% Section 2: Thermal Bremsstrahlung (``Free-Free Radiation'' From a Plasma)
% Merged from agent_01.tex through agent_11.tex
% Subsections 2.1--2.9

\section{Thermal Bremsstrahlung (``Free-Free Radiation'' From a Plasma)}
\label{sec:2}

\subsection{Thermal Bremsstrahlung (``Free-Free Radiation'' From a Plasma)}
\label{sec:2.1}

\textit{Basic references}\quad chapter~6 of Shkarofsky, Johnston, and
Bachynski~(1966)\cite{Shkarofsky1966};
chapter~15 of Jackson~(1962)\cite{Jackson1962}; \S4 of
Ginzburg~(1967)\cite{Ginzburg1967}.

\bigskip

\textit{Notation for fundamental constants}

\medskip

\noindent
\begin{tabular}{ll}
$h$            & Planck's constant \\
$m_e$          & rest mass of electron \\
$m_p$          & rest mass of proton \\
$e$            & charge of electron \\
$c$            & speed of light \\
$\alpha$       & fine-structure constant \\
$r_0$          & classical electron radius \\
$R_y$          & Rydberg energy \\
$Z$            & charge of ions in plasma \\
$\mathcal{A}$  & atomic weight of ions \\
$k$            & Boltzmann's constant \\
$\xi$          & Euler constant ($= 1.78\ldots$)
\end{tabular}

\bigskip

\textit{Radiation from a single classical collision}

\medskip

Consider a single ion of charge~$Z$, at rest in the laboratory frame. Let a single
electron of speed $v \ll c$ (nonrelativistic) scatter off the nucleus (Coulomb
scattering), with impact parameter~$b$. Let $I(\omega)$ be the total energy per unit
circular frequency, $\omega$, radiated by the electron as it scatters. A simple classical
calculation [e.g.\ Jackson~(1962)\cite{Jackson1962}, p.~507] gives
%
\begin{equation}
I(\omega) = \frac{2}{3} \frac{e^2}{m_e c^3} \left|\Delta \mathbf{v}_\omega\right|^2,
\tag{2.1.1}
\end{equation}
%
where $\Delta \mathbf{v}_\omega$ is the change in electron velocity during a time $\tau = 1/\omega$ centered
about the electron's point of closest approach to the nucleus.

It is useful to examine two limiting cases: small-angle scattering ($\theta \ll 1$)
corresponding to large impact parameter ($b \gg b_{s\text{-}l}$); and large-angle scattering
($\theta \sim 2\pi$), corresponding to small impact parameter ($b \ll b_{s\text{-}l}$). The impact
parameter, $b_{s\text{-}l}$, which separates small-angle from large-angle scattering, is given
by
%
\begin{equation}
Ze^2 / b_{s\text{-}l} = \tfrac{1}{2}\, m_e v^2.
\tag{2.1.2}
\end{equation}

\textit{For small-angle scattering} the total scattering angle $\theta$ is
%
\begin{equation}
\theta = \frac{1}{v}\, |\Delta \mathbf{v}|
      = \frac{Ze^2/b}{\frac{1}{2}\, m_e v^2}
      = \frac{b_{s\text{-}l}}{b}\,,
\tag{2.1.3}
\end{equation}
%
and the time during which the scattering occurs is $\Delta t \approx b/v$.  Hence, for
$\tau = 1/\omega \gg \Delta t$, the change in velocity during time $\tau$, $|\Delta \mathbf{v}_\tau|$,
is the full change $|\Delta \mathbf{v}|$ of~(2.1.3); while for $\tau = 1/\omega \ll \Delta t$ the change is
essentially zero.  Inserting these changes into equation~(2.1.1), we obtain
%
\begin{align}
I(\omega) &= 0  &\text{for}&\quad \omega \gg v/b, \notag \\[4pt]
I(\omega) &= \frac{8}{3}\,\frac{Z^2 e^2}{m_e c}\left(\frac{r_0}{v\,b}\right)^{\!2}
           &\text{for}&\quad \omega \ll v/b.
\tag{2.1.4}
\end{align}

\textit{For large-angle scattering} the electron orbit is very nearly a parabola, with
``perihelion'' at radius
%
\begin{equation}
r_p = b^2 / b_{s\text{-}l}
\tag{2.1.5a}
\end{equation}
%
and with speed at perihelion
%
\begin{equation}
v_p = c\bigl(2 Z r_0 / r_p\bigr)^{1/2}.
\tag{2.1.5b}
\end{equation}
%
During time intervals $\tau = 1/\omega \ll r_p/v_p$, a negligible amount of velocity change
occurs; hence
%
\begin{equation}
I(\omega) = 0 \qquad \text{for}\quad \omega \gg v_p/r_p
           = c(2Zr_0\, b_{s\text{-}l}^{\,3})^{1/2}/b^3.
\tag{2.1.6a}
\end{equation}

During time intervals $\tau = 1/\omega \gg r_p/v_p$ but $\tau < b/v$, the electron is initially
($\tau = -\tau/2$) headed directly toward the nucleus with ``parabolic'' speed

\noindent it swings into a sharp turn around the nucleus; and it emerges
at $t = +\tau/2$ with its initial velocity precisely reversed.
Consequently, $|\Delta v_\omega|$ is equal to~$2v_f$, and $I(\omega)$~is
\begin{equation}
  I(\omega)
  = \frac{32}{3}\,\frac{e^{2}}{c^{3}}
    \left(\frac{Zr_0\,\omega\,c^{2}}{3}\right)^{\!2/3}
  \quad\text{for}\quad
  \frac{v}{b_{s-1}} < \omega < \frac{v_p}{r_p}
  = \frac{c(2Zr_0\,b_{s-1}^{3})^{1/2}}{b^{3}}\,.
  \tag{2.1.6b}
\end{equation}

%=============================================================================
% Radiation from a nonrelativistic beam of electrons
%=============================================================================

\bigskip
\noindent\textit{Radiation from a nonrelativistic beam of electrons}

\medskip

Consider a single ion of charge~$Z$, at rest in the laboratory frame.
Bombard the ion with a monochromatic beam of electrons with speed
$v \ll 1$ and flux (number per unit time per unit area)~$S$.  Examine the
radiation of frequency~$\nu$ (circular frequency $\omega = 2\pi\nu$)
emitted by the electrons as they scatter off the nucleus.
Let $\mathscr{P}_\nu$ be the total power emitted per unit frequency
(ergs~sec$^{-1}$~Hz$^{-1}$); and define the emission cross section per
unit frequency, $d\sigma/d\nu$, by
\begin{equation}
  \mathscr{P}_\nu = (d\sigma/d\nu)\,S\,h\nu.
  \tag{2.1.7}
\end{equation}
The dependence of the emission cross section on photon frequency~$\nu$
and on electron speed~$v$ is most conveniently expressed in terms of
the two dimensionless parameters
\begin{equation}
  \frac{h\nu}{\tfrac{1}{2}m_e v^{2}}\,,
  \qquad
  \frac{\tfrac{1}{2}m_e v^{2}}{2R_y}
  = \left(\frac{v/c}{\alpha Z}\right)^{\!2}.
  \tag{2.1.8}
\end{equation}

On a sheet of paper (Figure~\ref{fig:2.1.1}) plot
$(h\nu)/(\tfrac{1}{2}m_e v^{2})$ vertically, and plot
$(\tfrac{1}{2}m_e v^{2})/(2R_y)$ horizontally.  This 2-dimensional plot
can be split into several different regions, in each of which the details
of the emission process are different.

%--- Forbidden region ---------------------------------------------------
\bigskip
\noindent\textit{Forbidden region}\quad
The individual photons emitted while scattering cannot have energies
greater than the electron kinetic energy---unless the electron gets
captured into a bound state around the ion.  But capture produces
``free-bound'' rather than bremsstrahlung radiation (\S\,2.2).
Consequently, the region
\begin{equation}
  (h\nu)/(\tfrac{1}{2}m_e v^{2}) > 1
  \tag{2.1.9a}
\end{equation}
of Figure~\ref{fig:2.1.1} is a ``forbidden region''; no bremsstrahlung
is emitted in this region; $d\sigma/d\nu$ vanishes in this region.

%--- Photon-discreteness region -----------------------------------------
\bigskip
\noindent\textit{Photon-discreteness region}\quad In the region
\begin{equation}
  \tfrac{1}{3} < (h\nu)/(\tfrac{1}{2}m_e v^{2}) \le 1
  \tag{2.1.9b}
\end{equation}
no more than 3~photons can be emitted by each scattering.  This ``photon
discreteness'' has a significant effect on the emission cross section,
cutting it down toward zero as one approaches the edge of the forbidden
region, $(h\nu)/(\tfrac{1}{2}m_e v^{2}) = 1$.  No classical calculation
can reveal such an effect; the photon-discreteness region must be
analyzed in a quantum mechanical manner; see, e.g., Heitler
\cite{Heitler1954}, and see below.

%=============================================================================
% Large-angle region (page 350)
%=============================================================================

\bigskip
\noindent\textit{Large-angle region}\quad
All small-angle scatterings last too long
($\Delta\tau \sim b/v > b_{s-1}/v$; cf.\ eq.~(2.1.4)) to produce any
radiation of $\tau = 1/\omega < b_{s-1}/v$; cf.\ eq.~(2.1.4).
Consequently, in the region
\begin{equation}
  \frac{h\nu}{\tfrac{1}{2}m_e v^{2}}
  > \frac{h\omega}{\tfrac{1}{2}m_e v^{2}}
  = \left(\frac{\tfrac{1}{2}m_e v^{2}}{Z^{2}R_y}\right)^{\!1/2}
  \tag{2.1.10}
\end{equation}
only large-angle scatterings produce radiation.  This region is shown in
Figure~\ref{fig:2.1.1}.  The power per unit frequency emitted at a fixed
frequency~$\nu$ in this region is
\begin{equation}
  \mathscr{P}_\nu
  = 2\pi \int_0^{b_{\max}} I(\omega)\,S\,2\pi\,b\,db\,,
  \tag{2.1.11a}
\end{equation}
where $I(\omega)$ is given by the large-angle formula~(2.1.6), and
$b_{\max}$ is the ``cutoff''
\begin{equation}
  b_{\max} = (2c^{2}\,Zr_0\,b_{s-1}^{3})^{1/6}
  \tag{2.1.11b}
\end{equation}
in that formula.

Performing the integration, dividing by $S\,h\nu$ to obtain
$d\sigma/d\nu$, and reexpressing the result in terms of fundamental
atomic constants, one obtains
\begin{equation}
  \frac{d\sigma}{d\nu}
  = \frac{2^{1/3}\,16\,\alpha\,c^{2}\,Z^{2}\,r_0^{2}}
         {3^{5/3}\;v^{2}}\,.
\end{equation}

A more exact classical calculation gives a slightly different numerical
coefficient:
\begin{equation}
  \left(\frac{d\sigma}{d\nu}\right)_{\!\text{LA}}
  = \frac{16\pi}{3\sqrt{3}}\;
    \frac{\alpha\,c^{2}\,Z^{2}\,r_0^{2}}{v^{2}}\,.
  \tag{2.1.12}
\end{equation}

\noindent(Here ``LA'' stands for ``large-angle region''.)

In other regions of Figure~\ref{fig:2.1.1}, one gets other formulas for
$d\sigma/d\nu$ (see below).  It is conventional in all regions to lump
the deviations from this large-angle cross section into a correction
$G(\nu,v)$ which is called the ``Gaunt factor'':
\begin{equation}
  G(\nu,v)
  \equiv \frac{(d\sigma/d\nu)}{(d\sigma/d\nu)_{\text{LA}}}\,.
  \tag{2.1.13}
\end{equation}

Thus, in the large-angle region $G = 1$; and in the forbidden region
$G = 0$.  Throughout the large-angle and photon-discreteness portion of
the large-angle region, $G$~remains approximately~1; see
Figure~\ref{fig:2.1.1}.

%=============================================================================
% Small-angle, classical region (pages 350--351)
%=============================================================================

\bigskip
\noindent\textit{Small-angle, classical region}\quad
In the region $\tau = 1/\omega \gg b_{s-1}/v$---i.e.
\begin{equation}
  \frac{h\nu}{\tfrac{1}{2}m_e v^{2}}
  \ll \left(\frac{\tfrac{1}{2}m_e v^{2}}{2R_y}\right)^{\!1/2},
  \tag{2.1.14}
\end{equation}
radiation is produced by small-angle scatterings as well as by
large-angle scatterings.  As a function of impact parameter~$b$, the
energy radiated by a single electron is constant in the large-angle
region, $b < b_{s-1}$ [eq.~(2.1.6b)], and decreases as~$1/b^{2}$ in
the small-angle region, $b > b_{s-1}$ [eq.~(2.1.4)].  Hence, when one
integrates over $2\pi\,b\,db$ to get the power radiated,
$\mathscr{P}_\nu$, one obtains

\smallskip
\noindent(contribution from small angles)
$\displaystyle{}= \ln\!\left(\frac{b_{\max}}{b_{s-1}}\right)$,

\smallskip
\noindent(contribution from large angles) ${}= \text{const}$,

\smallskip
\noindent where $b_{\max}$ is the cutoff in the small-angle region
(eq.~2.1.4):
\begin{equation}
  b_{\max} = v/\omega.
  \tag{2.1.15}
\end{equation}

Since the small-angle contribution is logarithmically dominant, one can
ignore the contribution from large-angle scatterings and calculate
\begin{equation}
  \frac{d\sigma}{d\nu}
  = \frac{\mathscr{P}_\nu}{S\,h\nu}
  = \frac{16\,\alpha\,c^{2}\,Z^{2}\,r_0^{2}}{3\,v^{2}}\,
    \ln\!\left(\frac{b_{\max}}{b_{s-1}}\right).
  \tag{2.1.16}
\end{equation}

\noindent A more exact classical calculation gives a slightly different
argument in the logarithm:
\begin{equation}
  \frac{d\sigma}{d\nu}
  = \frac{16\,\alpha\,c^{2}\,Z^{2}\,r_0^{2}}{3\,v^{2}}\,
    \ln\!\left(\frac{2\,b_{\max}}{\xi\;b_{s-1}}\right),
  \tag{2.1.18}
\end{equation}
where $\xi = 1.781\ldots$ is the ``Euler constant''.  Consequently, the
Gaunt factor for the small-angle, classical region is
\begin{equation}
  G(\nu,v)
  = \frac{\sqrt{3}}{\pi}\,
    \ln\!\left[
      \frac{2}{\xi}
      \left(\frac{\tfrac{1}{2}m_e v^{2}}{h\nu}\right)
      \left(\frac{\tfrac{1}{2}m_e v^{2}}{2R_y}\right)^{\!1/2}
    \right].
  \tag{2.1.19}
\end{equation}

Actually, this small-angle classical result is not valid throughout the
region~(2.1.14).  The uncertainty principle requires a modification at
electron energies $\tfrac{1}{2}m_e v^{2}$ larger than~$Z^{2}R_y$.

%=============================================================================
% Small-angle, uncertainty-principle region (pages 351--352)
%=============================================================================

\bigskip
\noindent\textit{Small-angle, uncertainty-principle region}\quad
Consider an electron with impact parameter~$b$ and speed~$v$.  The
electron is actually not a classical object; rather, it is a
quantum-mechanical wave packet.  To scatter with impact parameter~$b$,
the wave packet (i)~must have transverse dimensions, $\Delta x$, less
than~$b$
\begin{equation}
  \Delta x < b,
\end{equation}
and (ii)~must not spread transversely by more than~$b$ during the
classical scattering time $\Delta t = b/v$
\begin{equation}
  b > \left(\frac{\text{spreading}}{\text{during }\Delta t}\right)
  = \left(\frac{\Delta p_x}{m_e}\right)\Delta t
  > \left(\frac{\hbar}{m_e\,\Delta x}\right)\Delta t
  > \frac{\hbar}{m_e\,v}
  = \frac{b}{m_e\,v}\,\frac{\hbar}{b}\,.
\end{equation}

\noindent Thus, the uncertainty principle (``$\ge$'' in the above
equation) prevents the existence of scatterings with~$b$ less than the
electron de~Broglie wavelength
\begin{equation}
  b_{dB} = \hbar / m_e v.
  \tag{2.1.20}
\end{equation}

If $b_{dB} < b_{s-1}$, this limitation has no effect on small-angle
scatterings; and the small-angle, classical Gaunt factor~(2.1.19) is
valid.  If $b_{dB} > b_{s-1}$, the uncertainty principle comes into
play, and one must use~$b_{dB}$ rather than~$b_{s-1}$ as the lower
limit on the small-angle integral~(2.1.17) for $d\sigma/d\nu$.  The
result is
\begin{equation}
  \frac{d\sigma}{d\nu}
  = \frac{16\,\alpha\,c^{2}\,Z^{2}\,r_0^{2}}{3\,v^{2}}\,
    \ln\!\left(\frac{2\,b_{\max}}{\xi\;b_{dB}}\right),
  \tag{2.1.21}
\end{equation}
corresponding to the Gaunt factor
\begin{equation}
  G(\nu,v)
  = \frac{\sqrt{3}}{\pi}\,
    \ln\!\left(\frac{2\,b_{\max}}{\xi\;b_{dB}}\right)
  = \frac{\sqrt{3}}{\pi}\,
    \ln\!\left[
      \frac{2}{\xi}
      \left(\frac{\tfrac{1}{2}m_e v^{2}}{h\nu}\right)
    \right].
  \tag{2.1.22}
\end{equation}

The ``small-angle, uncertainty-principle region,'' in which this
expression holds, is separated from the ``small-angle, classical
region'' by the line
\begin{equation}
  b_{dB}/b_{s-1}
  = (\tfrac{1}{2}m_e v^{2})/(Z^{2}R_y) = 1.
\end{equation}

\noindent See Figure~\ref{fig:2.1.1}.

From Figure~\ref{fig:2.1.1} it is clear that the uncertainty principle
cannot have any effect on the large-angle region.

%=============================================================================
% Figure 2.1.1
%=============================================================================

\begin{figure}[htbp]
\centering
% \includegraphics[width=\columnwidth]{fig_2_1_1}
\fbox{\parbox{0.9\columnwidth}{\centering
[Figure 2.1.1 placeholder]\\[6pt]
Horizontal axis:
$\bigl(\tfrac{1}{2}m_e v^{2}\bigr)\big/\bigl(Z^{2}R_y\bigr)$.\quad
Vertical axis:
$h\nu\big/\bigl(\tfrac{1}{2}m_e v^{2}\bigr)$.\\[4pt]
Five regions shown:\\[2pt]
(1) Forbidden region ($G = 0$, $d\sigma/d\nu = 0$)
    --- above the line $h\nu/(\tfrac{1}{2}m_ev^2) = 1$.\\
(2) Photon-discreteness region ($G \lesssim 1$)
    --- between $\tfrac{1}{3}$ and~$1$.\\
(3) Large-angle region ($G = 1$)
    --- below discreteness, above
    $h\nu/(\tfrac{1}{2}m_ev^2) =
     [(\tfrac{1}{2}m_ev^2)/(2R_y)]^{1/2}$.\\
(4) Small-angle, classical region:
    $G = \dfrac{\sqrt{3}}{\pi}
    \ln\!\Bigl[\dfrac{4}{\xi}
    \bigl(\dfrac{\frac{1}{2}m_ev^{2}}{h\nu}\bigr)
    \bigl(\dfrac{\frac{1}{2}m_ev^{2}}{2R_y}\bigr)^{\!1/2}\Bigr]$
    --- lower-left.\\[4pt]
(5) Small-angle, U.P.\ (uncertainty-principle) region:
    $G = \dfrac{\sqrt{3}}{\pi}
    \ln\!\Bigl[\dfrac{4}{\xi}
    \bigl(\dfrac{\frac{1}{2}m_ev^{2}}{h\nu}\bigr)\Bigr]$
    --- lower-right.
}}
\caption{Regions and Gaunt factors for bremsstrahlung of frequency~$\nu$
emitted by an electron of kinetic energy~$\tfrac{1}{2}m_e v^{2}$
when it impinges on an ion of charge~$Ze$.}
\label{fig:2.1.1}
\end{figure}

out that throughout the photon-discreteness portion of the large-angle region,
$G$ remains approximately 1; see Figure~2.1.1.

\textit{Small-angle, classical region}\quad In the region $\tau = 1/\omega \gg b_{s-l}/v$---i.e.
\begin{equation}
\frac{h\nu}{\frac{1}{2}m_e v^2}
\ll \left(\frac{\frac{1}{2}m_e v^2}{Z^2 R_y}\right)^{1/2}
\tag{2.1.14}
\end{equation}

In the ``photon-discreteness'' portion of the small-angle, uncertainty-principle
region, discreteness effects modify the Gaunt factor (2.1.22) into the form
\begin{subequations}
\begin{align}
G(\nu,v) &= \frac{\sqrt{3}}{\pi}\,
\ln\!\left[\frac{\frac{1}{2}m_e(v_i + v_f)^2}{h\nu}\right],
\tag{2.1.24a}\\[6pt]
G(\nu,v) &= \frac{\sqrt{3}}{\pi}\,
\ln\!\left[\frac{v_i + v_f}{v_i - v_f}\right]^2,
\tag{2.1.24b}
\end{align}
\end{subequations}

Here $v_i$ and $v_f$ are the electron speeds before and after the collision:
\begin{equation}
\tfrac{1}{2}m_e v_f^2 = \tfrac{1}{2}m_e v_i^2 - h\nu.
\end{equation}

[Expression (2.1.24) is the result of a quantum-mechanical Born-approximation
calculation, valid throughout the small-angle, uncertainty-principle region.]

%%---------------------------------------------------------------------------
%% Radiation from a plasma
%%---------------------------------------------------------------------------

\bigskip
\textit{Radiation from a plasma}

Consider a plasma which may contain a variety of ionic species.  Focus attention
on radiation from all ions of a particular type, with charge $Ze$ and atomic weight
$A$.  Let $C_A$ be the concentration by mass of the species of atom that gives rise to
this ion, and let $f_i$ be the fraction of all such atoms ionized to the state of interest.
Let $f_e$ be the number of unbound electrons per baryon in the gas.  Then
\begin{equation}
n_I = \text{(number of ions per unit volume)} = f_i C_A \rho_0 / A m_p.
\tag{2.1.25}
\end{equation}

Here $\rho_0$ is the density of rest mass.  (Of course, $f_i \leq 1$; $C_A \leq 1$.)  The ions
and electrons both have Maxwell velocity distributions; but because $m_e \ll A m_p$,
the velocities of the electrons and their accelerations during Coulomb scattering
are far greater than those of the ions.  Therefore, in calculating bremsstrahlung
one can regard the ions as at rest.  The number of photons per unit mass of plasma
is $f_e/m_p$, so the total power per unit frequency emitted from one gram of plasma
by electron-ion scatterings is
\begin{equation}
\varepsilon_\nu
= \frac{C_A f_i}{A m_p}
\underbrace{\vphantom{\bigg|}\quad}_{\substack{
\text{fractions of all} \\
\text{ions per} \\
\text{unit mass}
}}
\cdot \int
\underbrace{\vphantom{\bigg|}\quad}_{\substack{
\text{number} \\
\text{of electrons}
}}
\underbrace{\vphantom{\bigg|}\quad}_{\substack{
\text{electrons that have} \\
\text{speeds $v$ in range $dv$}
}}
\frac{f_e \rho_0}{m_p}\;v\,h\nu\;
\frac{d\sigma}{dv}\;dv.
\tag{2.1.26}
\end{equation}

Here $T$ is the temperature in the plasma.  By inserting into the integrand
expressions (2.1.12) and (2.1.13) for $d\sigma/dv$ and performing the integration, one
obtains
\begin{align}
\varepsilon_\nu
&= \frac{C_A f_i f_e^2}{A m_p}\,
\frac{32\pi e^6}{3\,m_e c^3}\,
\left(\frac{2\pi m_e}{3kT}\right)^{\!1/2}\!
\frac{\rho_0^2}{m_p^2}\;
e^{-h\nu/kT}\;
\bar{G}(\nu,T)
\notag\\[4pt]
&= 2.5 \times 10^{10}\;
\frac{\text{erg}}{\text{g\;sec\;Hz}}\;
\left(\frac{C_A f_i f_e^2}{A}\right)
\left(\frac{\rho_0}{\text{g\;cm}^{-3}}\right)
\left(\frac{\rho_0}{A}\right)
T_K^{\,1/2}\;e^{-h\nu/kT}\;\bar{G}(\nu,T).
\tag{2.1.27}
\end{align}

%%---------------------------------------------------------------------------
%% Mean Gaunt factor \bar{G}
%%---------------------------------------------------------------------------

Here $T_K$ is temperature measured in ${}^\circ$K, and $\bar{G}(\nu,T)$ is the ``Maxwell--Boltzmann-averaged
Gaunt factor'', also called the ``mean Gaunt factor'':
\begin{equation}
\bar{G}(\nu,T) = \int_0^\infty
G\!\left(\nu,\;v = \left[\frac{2(xkT + h\nu)}{m_e}\right]^{1/2}\right)
e^{-x}\,dx.
\tag{2.1.28}
\end{equation}

Corresponding to the 2-dimensional plot of Gaunt factors (Figure~2.1.1) one
can make a 2-dimensional plot of mean Gaunt factors (Figure~2.1.2).  It is
straightforward to calculate the mean Gaunt factors shown there (up to
quantities of order unity) from the Gaunt factors of Figure~2.1.1.  Notice that,

%%---------------------------------------------------------------------------
%% Figure 2.1.2 — Mean Gaunt factors
%%---------------------------------------------------------------------------

\begin{figure}[htbp]
\centering
% \includegraphics[width=\columnwidth]{fig_2_1_2}
\fbox{\parbox{0.9\columnwidth}{\centering [Figure 2.1.2 placeholder]\\[6pt]
Mean Gaunt factors for bremsstrahlung of frequency $\nu$ emitted by
electron-ion collisions in a plasma of temperature $T$.\\[4pt]
Horizontal axis: $h\nu/kT$.\\
Vertical axis: $(kT)/(Z^2 R_y)$.\\[4pt]
Regions shown:\\
--- ``Small-angle, u.p., tail region'',
$\bar{G} = \bigl(\frac{3}{\pi}\bigr)\,\bigl(\frac{kT}{h\nu}\bigr)^{1/2}$\\
--- ``Small-angle, u.p.\ region'',
$\bar{G} = \frac{\sqrt{3}}{\pi}\ln\!\left[\frac{4}{\xi}\,\frac{kT}{h\nu}\right]$\\
--- ``Small-angle, classical region'',
$\bar{G} = \frac{\sqrt{3}}{\pi}\,\ln\!\left[\left(\frac{4}{\xi^{3/2}}\right)
\left(\frac{kT}{h\nu}\right)\left(\frac{kT}{Z^2 R_y}\right)^{1/2}\right]$\\
--- ``Large-angle region'', $\bar{G} = 1$\\
--- ``Large-angle, tail region'', $\bar{G} = 1$\\
--- ``Forbidden region''}}
\caption{\textit{Figure~2.1.2.}\ Mean Gaunt factors for bremsstrahlung of frequency $\nu$ emitted by
electron-ion collisions in a plasma of temperature~$T$.}
\label{fig:2.1.2}
\end{figure}

aside from quantities of order unity, the mean Gaunt factors in the large-angle
and small-angle regions are obtained by simply taking the Gaunt
factors and making the replacement $\frac{1}{2}m_e v^2 \approx kT$:
\begin{equation}
\bar{G}(\nu,kT) \simeq G\bigl(\nu,\;\tfrac{1}{2}m_e v^2 = kT\bigr).
\end{equation}

Notice also that the ``forbidden region'' of the Gaunt-factor plot is replaced, in
the mean-Gaunt-factor plot, by ``tail regions''.  The radiation in these tail regions
is produced by the high-energy tail, $\frac{1}{2}m_e v^2 \approx h\nu \gg kT$, of the Maxwell--Boltzmann
distribution.

Tables of mean Gaunt factors, calculated with much higher accuracy than
Figure~2.1.2, are given by \citet{Green1959} and by \citet{Karzas1961}.

From Figure~2.1.1 it is clear that the uncertainty principle cannot have any
effect on the large-angle region.

%%---------------------------------------------------------------------------
%% Total emissivity and frequency-averaged mean Gaunt factor
%%---------------------------------------------------------------------------

The total energy emitted by electron-ion collisions is obtained by integrating
expression (2.1.27) over frequency:
\begin{align}
e &= \int \varepsilon_\nu\,d\nu
= \frac{16\,C_A f_i f_e^2}{3}\;
\frac{\alpha\,m_e c^2_0}{A}\;
\left(\frac{2\pi kT}{m_e c^2}\right)^{\!1/2}\!
\frac{\rho_0}{m_p^2}\;
\left(\frac{C_A f_i f_e^2}{A}\right)
\left(\frac{\rho_0}{\text{g\;cm}^{-3}}\right)
T_K^{\,1/2}\;\bar{G}(T)
\notag\\[4pt]
&= 5.2 \times 10^{20}\;
\frac{\text{ergs}}{\text{g\;sec}}\;
\left(\frac{C_A f_i f_e^2}{A}\right)
\left(\frac{\rho_0}{\text{g\;cm}^{-3}}\right)
T_K^{\,1/2}\;\bar{G}(T).
\tag{2.1.29}
\end{align}

Here $\bar{G}(T)$ is the frequency-averaged mean Gaunt factor
\begin{equation}
\bar{G}(T) = \int_0^\infty \bar{G}(\nu = xkT/h,\;T)\;e^{-x}\,dx.
\tag{2.1.30}
\end{equation}

From a value of $\bar{G} \simeq 1.2$ at $kT/Z^2 R_y \sim 0.01$, it increases gradually to a maximum
of $1.42$ at $kT/Z^2 R_y \sim 1$, and then decreases gradually toward an asymptotic,
high-temperature value of $1.103$ at $kT/Z^2 R_y > 100$; see Figure~2 of \citet{Green1959}.

In a plasma at nonrelativistic temperatures, the radiation from electron-electron
collisions and from ion-ion collisions is negligible compared to electron-ion
radiation.  This is because (i) when identical particles scatter, the first time
derivative of the electric dipole moment is conserved, so only radiation of
quadrupole order and higher [with intensity $\sim(v/c)^2$ less than dipole radiation]
can be emitted; (ii) the speeds and accelerations of colliding ions are far less than
those of electrons because of their much larger masses.

Equations (2.1.27) and (2.1.29), together with Figures~2.1.2 and~2.1.3, are
the chief results of astrophysical interest from this section.  When dealing with
atoms that are only partially ionized, one must sometimes correct these results
for screening of the nuclear charge by the bound electrons; see, e.g.,
\citet{Ginzburg1967}.  For temperatures approaching and exceeding $kT = m_e c^2$,
relativistic effects and electron-electron collisions modify the emissivity
(2.1.29) to the form
\begin{align}
e &= 5.2 \times 10^{20}\;
\frac{\text{ergs}}{\text{g\;sec}}\;
\left(\frac{C_A f_i f_e^2}{A}\right)
\left(\frac{\rho_0}{\text{g\;cm}^{-3}}\right)
T_K^{\,1/2}\bigl(1 + 4.4 \times 10^{-10}\,T_K\bigr)\,\bar{G}(T).
\tag{2.1.31}
\end{align}

\noindent (\citealt{Ginzburg1967}).


%%---------------------------------------------------------------------------
%%  SECTION 2.2 — Free-Bound Radiation
%%---------------------------------------------------------------------------

\subsection{Free-Bound Radiation}
\label{sec:2.2}

\textit{Basic references}\quad \citet{Jackson1962}; \citet{Ginzburg1967}.

\bigskip
\textit{Basic physics and formulas}

When an electron of kinetic energy
\begin{equation}
\tfrac{1}{2}m_e v^2 \leq Z^2 R_y
\end{equation}
impinges on an ion, there is a significant probability that, instead of scattering
with the emission of bremsstrahlung, it will get captured into a bound state with
an accompanying emission of ``free-bound'' radiation.

We shall not compute here the details of the radiation.  Instead we cite the
results of such a computation from the above references.

Consider a plasma at temperature $T$ and density of rest mass $\rho_0$.  Let $C_A$ be
the concentration by mass of a particular atom in the plasma; let $f_i$ be the fraction
of all such atoms ionized into a given state, with charge $Ze$; and let $n_q$ be the
principal quantum number of that state, and $E_i$ its
ionization energy.  Then the emissivity due to such transitions is
\begin{align}
\varepsilon_\nu &= 8c^2
\left(\frac{C_A f_i f_e^2}{An_q^5}\right)
\left(\frac{\alpha^2 m_e e^2}{m_p}\right)^{\!3/2}
\rho_0
\left(\frac{\rho_0}{3kT}\right)^{\!3/2}
e^{-(h\nu - E_i)/kT}\;G_{bf}
\notag\\[4pt]
&= 3.9 \times 10^{15}\;
\frac{\text{ergs}}{\text{g\;sec\;Hz}}\;
\left(\frac{C_A f_i f_e^2}{An_q^5}\right)
\left(\frac{\rho_0}{\text{g\;cm}^{-3}}\right)
T_K^{-3/2}\;e^{-(h\nu - E_i)/kT}\;G_{bf}.
\tag{2.2.1}
\end{align}

Here $G_{bf}$ is a ``bound-free Gaunt factor'', analogous to the ``free-free'' Gaunt
factors of \S\ref{sec:2.1}.  It depends on the kinetic energy
\begin{equation}
\tfrac{1}{2}m_e v^2 = h\nu - E_i
\tag{2.2.2}
\end{equation}
of the electron that is captured, and on the structure of the bound state (which
one characterizes by its quantum numbers $n_q, l_q, \ldots$):
\begin{equation}
G_{bf} = G_{bf}(h\nu - E_i;\;n_q, l_q, \ldots).
\tag{2.2.3}
\end{equation}

Gaunt factors for various bound-free transitions of astrophysical interest are
tabulated and graphed by \citet{Karzas1961}.  Energy conservation
requires that all photons emitted have energies greater than $E_i$; consequently,
\begin{equation}
G_{bf} = 0 \quad \text{for} \quad h\nu - E_i < 0.
\tag{2.2.4}
\end{equation}

In general, $G_{bf}$ ``turns on'' discontinuously at $h\nu = E_i$, with an initial value
between ${\sim}\,0.5$ and ${\sim}\,1.5$:
\begin{equation}
G_{bf} \simeq 1 \quad \text{for} \quad 0 < h\nu - E_i \lesssim E_i.
\tag{2.2.5}
\end{equation}

For $h\nu - E_i \gtrsim Z^2 R_y$, $G_{bf}$ typically remains within an order of magnitude of
unity; as $h\nu - E_i \geq 10\,Z^2 R_y$, $G_{bf}$ falls rapidly.
See \citet{Karzas1961}.

The total emissivity, $e = \int \varepsilon_\nu\,d\nu$, due to a given free-bound transition is
\begin{align}
e &= 8\sqrt{3}\,c^2
\left(\frac{C_A f_i f_e^2}{An_q^5}\right)
\left(\frac{\alpha^2 m_e e^2}{m_p^2}\right)^{\!3/2}
\rho_0
\left(\frac{2\pi m_e c^2}{3kT}\right)^{\!1/2}
T_K^{-1/2}\;\bar{G}_{bf}
\notag\\[4pt]
&= 4.2 \times 10^{26}\;
\frac{\text{ergs}}{\text{g\;sec}}\;
\left(\frac{C_A f_i f_e^2}{An_q^5}\right)
\left(\frac{\rho_0}{\text{g\;cm}^{-3}}\right)
T_K^{-1/2}\;\bar{G}_{bf}
\tag{2.2.6}
\end{align}

where $\bar{G}_{bf}$, the mean Gaunt factor, depends only on temperature and on the
bound state, and is approximately equal to one.
\begin{equation}
\bar{G}_{bf} = \int_0^\infty G_{bf}(h\nu - E_i = xkT;\;n_q, l_q, \ldots)\;e^{-x}\,dx \simeq 1.
\tag{2.2.7}
\end{equation}

The total emissivity (2.2.6) due to a given free-bound transition, compared
to the total emissivity (2.1.29) of free-free radiation from the same type of ion,
is
\begin{equation}
\frac{\varepsilon_{fb}}{\varepsilon_{ff}}
\simeq \frac{Z^2}{n_q^5}\,
\left(\frac{8 \times 10^5}{T}\right).
\tag{2.2.8}
\end{equation}

\citet{Ginzburg1967} states that for astrophysical plasmas the ratio of total free-bound
radiation to total free-free radiation (summed over all states and all ions)
generally is less than
\begin{equation}
\frac{\varepsilon_{fb\,\text{total}}}{\varepsilon_{ff\,\text{total}}}
< \left(\frac{8 \times 10^5\,\text{K}}{T}\right).
\tag{2.2.9}
\end{equation}

\subsection{Thermal Cyclotron and Synchrotron Radiation}
\label{sec:2.3}

\textit{Basic references}\quad Jackson~\cite{Jackson1962};\
Ginzburg~\cite{Ginzburg1967}.

\bigskip
\noindent\textit{Power radiated by a single electron}

\medskip
Consider an electron of speed~$v$ and total mass-energy~$\gamma m_e c^2$,
\begin{equation}
\gamma \equiv (1 - v^2)^{-1/2},
\tag{2.3.1}
\end{equation}
spiralling in a magnetic field of strength~$B$. The Lorentz 4-acceleration of the
electron is
\begin{equation}
a^0 = 0, \qquad
\mathbf{a} = (e / m_e c)\,\gamma\,\mathbf{v} \times \mathbf{B}.
\tag{2.3.2}
\end{equation}
Consequently, the total power radiated---as obtained from the standard
Lorentz-invariant equation,
\begin{equation}
\frac{dE}{dt} = \frac{2e^2}{3c^3}\,a^2
\tag{2.3.3}
\end{equation}

\noindent---is
\begin{equation}
\frac{dE}{dt} = \frac{2\,r_0^2}{3\,c}\,(\gamma v_\perp)^2\,B^2.
\tag{2.3.4}
\end{equation}
Here $v_\perp$ is the component of the electron velocity perpendicular to the magnetic
field. If the electron is nonrelativistic, $v \ll c$, this radiation is called ``cyclotron
radiation''. If the electron is ultrarelativistic, $\gamma \gg 1$, it is called ``synchrotron
radiation''.

\bigskip
\noindent\textit{Cyclotron radiation from a nonrelativistic plasma}

\medskip
Consider a plasma with temperature
\begin{equation}
kT \ll m_e c^2, \quad \text{i.e.}\ T \lesssim 6 \times 10^9\;\text{K},
\tag{2.3.5}
\end{equation}
and with a magnetic field of strength~$B$. As in the nonrelativistic case, radiation
from protons spiralling
in the magnetic field is smaller by a factor $(m_e/m_p)^3 \simeq 10^{-10}$ than that from
electrons, since
\begin{equation}
\frac{dE}{dt} \propto \frac{v^2}{m^2} \propto \frac{1}{m^3}\,.
\tag{2.3.6}
\end{equation}
(Recall: in thermal equilibrium the mean kinetic energies of protons and electrons
are the same.) Hence, we can ignore protons and other ions when calculating
cyclotron radiation. Let $f_e$ be the number of unbound electrons per baryon in
the plasma. Then the total emissivity (ergs per second per gram of plasma) is
\begin{equation}
\varepsilon = f_e \left\langle \frac{dE}{dt} \right\rangle
\bigg/\!\left(\frac{1}{m_p}\right),
\tag{2.3.7}
\end{equation}
where $\langle\;\rangle$ denotes an average over the Maxwell-Boltzmann velocity distribution
of the electrons. Since $\gamma = 1$, since 2 of the 3 spatial directions are orthogonal
to~$B$, and since the mean kinetic energy per electron is $\tfrac{3}{2}kT$, we have
\begin{equation}
\langle \gamma v_\perp^2 \rangle
= \frac{2}{3}\,\frac{2\,\bar{E}}{m_e}
= \frac{2}{3}\,\frac{3kT}{m_e}
= \frac{2kT}{m_e}\,.
\tag{2.3.8}
\end{equation}
Combining this with equations~(2.3.4) and~(2.3.7), we obtain for the total
emissivity of the plasma
\begin{align}
\varepsilon &= \frac{4}{3}\left(\frac{f_e\,r_0^2}{m_p}\right)
\left(\frac{kT}{m_e c^2}\right) B^2
\tag{2.3.9}\\
&= (0.32\;\text{ergs/g\,sec})\;f_e\;T_{\rm K}\;B_G^2\,,
\notag
\end{align}
where $T_{\rm K}$ is temperature in~${}^\circ$K and $B_G$ is magnetic field in Gauss.

Each electron emits its radiation monochromatically, with the cyclotron
frequency (orbital frequency of electron's spiral motion)
\begin{equation}
\nu_{\text{cyc}} = \frac{eB}{2\pi m_e c}
= (2.79\;\text{MHz})\;B_G\,.
\tag{2.3.10}
\end{equation}

Thus, if the magnetic field is uniform, the radiation will be monochromatic
with this frequency. But in Nature the magnetic field will always be inhomo\-geneous,
and the spectrum will show a spread proportional to the spread in~$B^2$.
(Of course, there are always other sources of spread in the spectrum---e.g., doppler
shifts and relativistic generation of harmonics.)

\textit{Synchrotron radiation from a relativistic plasma.}

\medskip
\noindent
Consider a plasma with temperature
\begin{equation}
kT \gg m_e c^2, \quad \text{i.e.}\ T \gg 6 \times 10^9\,\text{K},
\tag{2.3.11}
\end{equation}
and with a magnetic field of strength $B$.
As in the nonrelativistic case, radiation
from protons and ions is totally negligible.
The total emissivity due to electrons
and positrons (recall: at $kT \gg m_e c^2$, there will be many electron-positron pairs!)
is given, as before, by equation~(2.3.7); but now the factor $f_e$ must be the number
of electrons and positrons per baryon,\footnote{The case of an optically thin plasma
is of particular interest in astrophysics. The kinetic theory of the creation and
annihilation of positrons in this case is somewhat complex; see, e.g.,
\citet{BisnovatyiKogan1971}.}
and the relevant, ultrarelativistic Maxwell--Boltzmann average is

\begin{equation}
\langle (\gamma v_\perp)^2 \rangle = \tfrac{2}{3}\langle\gamma^2\rangle
= 8\left(\frac{kT}{m_e c^2}\right)^{\!2}.
\tag{2.3.12}
\end{equation}
[One evaluates $\langle\gamma^2\rangle$ using the distribution function in phase space
\begin{equation}
\mathscr{N} = \frac{dN}{d^3 x\,d^3 p} \propto e^{-E/kT} = e^{-\gamma m_e c^2/kT},
\tag*{}
\end{equation}
with
\begin{equation}
d^3 p = 4\pi c^{-3}E^2\,dE
\tag*{}
\end{equation}
in ultrarelativistic limit;
in particular
\begin{equation}
\langle\gamma^2\rangle = \frac{1}{(m_e c^2)^2}\,
\frac{\displaystyle\int_0^\infty E^2\,e^{-E/kT}\,E^2\,dE}
{\displaystyle\int_0^\infty e^{-E/kT}\,E^2\,dE}
= 12\left(\frac{kT}{m_e c^2}\right)^{\!2}.\,]
\tag*{}
\end{equation}

Combining equations~(2.3.12), (2.3.7), and~(2.3.4), we obtain for the total
emissivity of the plasma
\begin{align}
\varepsilon &= \frac{16}{3}\left(\frac{r_0^2 c}{m_p}\right)
\left(\frac{kT}{m_e c^2}\right)^{\!2} B^2 \notag \\
&= (2.2 \times 10^{-10}~\text{ergs/g\,sec})\,f_e\,T_K^2\,B_G^2.
\tag{2.3.13}
\end{align}

In the ultrarelativistic case the electrons do not emit monochromatically.
An electron of energy $\gamma m_e c^2$ beams most of its radiation into a forward cone of
half-angle $\alpha = 1/\gamma$ (special-relativistic ``headlight effect'').
Since the electron is
headed toward the observer with speed $v$ when its cone is directed toward him,
the cone sweeps past the observer in time
\begin{equation}
\Delta t = (1-v^2)\left(\frac{2\alpha}{\omega_{\text{cyc}}}\right) = (1-v^2)\,\frac{1}{\gamma^3\,\omega_{\text{cyc}}}
= \frac{1}{\gamma^3\,\omega_{\text{cyc}}}.
\tag{2.3.14}
\end{equation}
Here $\omega_{\text{cyc}}$ is the angular velocity of the electron in its spiralling orbit
\begin{equation}
\omega_{\text{cyc}} = \frac{eB}{\gamma m_e c}.
\tag{2.3.15}
\end{equation}

Thus, the radiation comes in short bursts, of duration $\sim\!\Delta t = 1/(\gamma^3\omega_{\text{cyc}})$,
separated by long intervals of time $2\pi/\omega_{\text{cyc}}$.
When Fourier analyzed, such
radiation must be concentrated near the ``critical frequency''
\begin{equation}
\nu_{\text{crit}} = \tfrac{3}{2}\,\nu_{\text{cyc}} = \gamma^2\,eB/m_e c.
\tag{2.3.16}
\end{equation}

A detailed calculation [chapter~14 of \citet{Jackson1962}] reveals a fairly broad
spectrum which rises as $\nu^{1/3}$ at low frequencies, which peaks at
\begin{equation}
\nu_{\text{peak}} = 0.29\,\nu_{\text{crit}},
\tag{2.3.17}
\end{equation}
and which decays exponentially $\sim\nu^{1/2}\,e^{-2\nu/\nu_{\text{crit}}}$, at $\nu \gg \nu_{\text{crit}}$.
When averaged over the relativistic plasma, such a spectrum will have a broad
peak, with maximum located at
\begin{equation}
\nu_{\text{peak, plasma}} = \frac{\gamma_{\text{peak}}^2\,eB}{4\,m_e c}
\simeq 12(kT/m_e c^2)^2
\quad\left[\text{value of $\gamma_{\text{peak}}^2$ for electrons of}\right]
\tag{2.3.18}
\end{equation}
\begin{equation}
\nu_{\text{peak, plasma}} \simeq (100~\text{MHz})\,T_{10}^2 B_G.
\tag*{}
\end{equation}
Here $T_{10}$ is temperature in units of $10^{10}$\,K.
At $\nu \ll \nu_{\text{peak, plasma}}$
the spectrum will rise as a power law.

%%% ====================================================================
%%% SECTION 2.4 — Electron Scattering of Radiation
%%% ====================================================================

\subsection{Electron Scattering of Radiation}
\label{sec:2.4}

\textit{Basic references}

\citet{Jackson1962}, \S\,14.7 of \citet{Jackson1962b};
\S\,12.3 of \citet{Leighton1959};
\citet{Kompaneets1957}; \citet{Weyman1965};
\citet{Sunyaev1973}.

\medskip
\noindent
\textit{Basic physics and formulas}

\medskip
\noindent
In discussing electron scattering, we shall confine attention to the nonrelativistic
case---i.e.\ we shall demand that the temperature of the gas and the photon
frequencies satisfy
\begin{equation}
kT \ll m_e c^2 \quad (T \ll 6 \times 10^9\,\text{K}); \quad
h\nu \ll m_e c^2 = 500~\text{keV}.
\tag{2.4.1}
\end{equation}

The differential cross section for a free electron to scatter a photon has the
form
\begin{equation}
\frac{d\sigma_e}{d\Omega} = \tfrac{1}{2}\,r_0^2\,(1 + \cos^2\theta),
\tag{2.4.2}
\end{equation}
where $\theta$ is the angle between incoming and outgoing photon.  The total
cross section (``Thomson cross section''), obtained by integrating over all
directions, is
\begin{equation}
\sigma_T = (8\pi/3)\,r_0^2 = 0.665 \times 10^{-24}~\text{cm}^2.
\tag{2.4.3}
\end{equation}
It is important for astrophysical applications that the scattering cross section
is ``color blind'' (no dependence on frequency).  See, e.g., \S\,4.

It is also important that the differential cross section is symmetric between
forward and backward directions.  This guarantees, for example, that on the
average a scattered photon transmits \textit{all} of its momentum to the electron
(Here $\mathbf{n}_s$ is a unit vector in the direction of the initial photon motion).  By
contrast, only a tiny fraction of the photon energy is transmitted to the electron.
In a frame where the electron is initially at rest, it receives a kinetic energy
\begin{subequations}
\begin{equation}
\Delta E_\gamma = (h\nu/m_e c^2)\,h\nu\,(1-\cos\theta).
\tag{2.4.4a}
\end{equation}
\begin{equation}
(\Delta\mathbf{p}_e) = \mathbf{p}_\gamma = (h\nu/c)\,\mathbf{n}_s.
\tag{2.4.4b}
\end{equation}
\end{subequations}
(Recall: we have assumed $h\nu/m_e c^2 \ll 1$.)

Consider a monochromatic beam of photons, with frequency $\nu$, moving
through a thermalized plasma of temperature $T$.  Scattering by protons and ions
will be negligible compared to scattering by electrons, since
\begin{equation}
\sigma_e \propto r_0^2 = (e^2/m_e c^2)^2 \propto 1/(\text{mass of scatterer})^2.
\tag{2.4.5}
\end{equation}

On the average, how much energy $\langle\Delta E_\gamma\rangle$ is transmitted to the electron in a
single scattering?  The answer for electrons nearly at rest ($kT \ll h\nu$) is obtained
by averaging expression~(2.4.4b) over all angles.  Because of the forward-backward
symmetry of the differential cross section, $\cos\theta$ averages to zero and one obtains
\begin{equation}
\langle\Delta E_\gamma\rangle = (h\nu/m_e c^2)\,h\nu, \quad \text{if } kT \ll h\nu.
\tag{2.4.6}
\end{equation}

One can calculate $\langle\Delta E_\gamma\rangle$ for higher temperatures by invoking conservation
of 4-momentum, and averaging over all angles and electron speeds.  However,
such a calculation is rather long and messy.  To obtain the answer more easily,
one can use the following trick: An examination of the law of 4-momentum
conservation convinces one that $\langle\Delta E_\gamma\rangle$ must have the temperature dependence
\begin{equation}
\langle\Delta E_\gamma\rangle \propto (h\nu - \alpha kT),
\tag{2.4.7}
\end{equation}
where $\alpha$ is a constant to be calculated.  (Thus, if $h\nu > \alpha kT$, the electrons get
heated by the photons; if $h\nu < \alpha kT$, they get cooled.)
This law will reduce to
expression~(2.4.6) for $kT \ll h\nu$ if and only if the proportionality constant is
$h\nu/m_e c^2$:
\begin{equation}
\langle\Delta E_\gamma\rangle = (h\nu/m_e c^2)(h\nu - \alpha kT).
\tag*{}
\end{equation}

To calculate the constant $\alpha$, imagine the following experiment.  Place a large
number of photons and electrons into a box with perfectly reflective walls.
Require that the photons and electrons interact only by electron scattering.  Then
photons cannot be created or destroyed; only scattered.  Hence, when thermal
equilibrium is reached, the photons acquire a Boltzmann energy distribution,
rather than a Planck distribution.  (Recall that stimulated emissions are responsible
for Planckian deviations from the Boltzmann law; \S\,2.5.)  Since photons have
zero rest mass, their Boltzmann distribution
\begin{equation}
\mathscr{N} = \frac{dN}{d^3 x\,d^3 p} \propto e^{-h\nu/kT},
\tag*{}
\end{equation}
corresponds to a number of photons per unit frequency given by
\begin{equation}
\frac{dN}{d\nu} \propto \nu^2\,e^{-h\nu/kT},
\tag{2.4.8}
\end{equation}
and corresponds to
\begin{equation}
\langle h\nu\rangle = 3kT, \qquad
\langle (h\nu)^2\rangle = 12(kT)^2.
\tag*{}
\end{equation}
Thus, for our thought experiment the energy transfer in each collision, expression
(2.4.7) averaged over the equilibrium photon distribution, is
\begin{equation}
\langle\langle\Delta E_\gamma\rangle\rangle = (3kT/m_e c^2)(4-\alpha)\,kT.
\tag*{}
\end{equation}

But in equilibrium there must be no average energy transfer between photons
and electrons; $\langle\langle\Delta E_\gamma\rangle\rangle$ must be zero.
This condition tells us that $\alpha = 4$.
Thus, turning back to the general situation, we conclude that monochromatic
photons passing through a plasma of temperature $T$ transfer an average energy
per collision
\begin{equation}
\langle\Delta E_\gamma\rangle = (h\nu/m_e c^2)(h\nu - 4kT)
\tag{2.4.9}
\end{equation}
to the electrons.  When $4kT \gg h\nu$, the photon energies get boosted by the
collisions, and one says that the radiation is being ``Comptonized''; see \S\,5.10;
see, e.g., \citet{Illarionov1972}.

This concludes, for the moment, our discussion of the interaction between
radiation and plasmas.  We must point out that we have ignored a number of
plasma effects and instabilities which can be important in astrophysical
situations.  See, e.g., \citet{Bekefi1966} and \citet{KaplanTsytovich1972}.

%%% ====================================================================
%%% SECTION 2.5 — Hydrodynamics and Thermodynamics
%%% ====================================================================

\subsection{Hydrodynamics and Thermodynamics}
\label{sec:2.5}

\textit{Basic references}\quad
\S\S\,22.2 and 22.3 of \citet{MTW1973};
\citet{Lichnerowicz1967} and references cited therein.
We adopt the notational conventions of \citet{MTW1973},
including signature ``$-+++$''; geometrized units ($c = G = 1$); Greek
indices ranging from 0 to 3 and Latin from 1 to 3; ``hats'' on indices that refer
to local orthonormal frames, e.g.\ $\mu^{\hat{\alpha}}$ and $T^{\hat{\alpha}\hat{\beta}}$; extra-bold, sans-serif type for
4-vectors and 4-tensors, e.g.\ $\fvec{u}$ and $\mathbf{T}$; normal bold-face type for 3-vectors and
3-tensors (local Euclidean geometry).

\medskip
\noindent
\textit{Parameters describing the ``fluid.''}

\medskip
\noindent
``LRF'' $=$ local rest frame of the baryons; i.e.\ local orthonormal frame in which
there is no net baryon flux in any direction.

\begin{itemize}
\item[$\fvec{u}$] 4-velocity of LRF; i.e.\ 4-velocity of the ``fluid''.
\item[$n$] number density of baryons (number of baryons per unit volume) as
measured in LRF.
\item[$m_B$] mean rest mass of a baryon in the fluid; i.e.
\end{itemize}
\begin{equation}
m_B = \left(\!\!
\begin{array}{c}
\text{rest mass of}\\
\text{hydrogen atom}
\end{array}\!\!\right)
\times
\left(\!\!
\begin{array}{c}
\text{fraction of all baryons that are}\\
\text{in the form of hydrogen nuclei}
\end{array}\!\!\right)
+
\frac{1}{4}
\left(\!\!
\begin{array}{c}
\text{rest mass of}\\
\text{helium atom}
\end{array}\!\!\right)
\times
\left(\!\!
\begin{array}{c}
\text{fraction of all baryons that}\\
\text{are in helium nuclei}
\end{array}\!\!\right)
+ \cdots
\tag{2.5.1}
\end{equation}

\textit{Note}: in this section we restrict ourselves to fluids that are chemically
homogeneous and in which no nuclear reactions occur (``standard
fluid'').  Thus, $m_B$ is constant.

\begin{align}
\rho_0 &\equiv m_B\,n, && \text{rest-mass density, defined by}
\tag{2.5.2} \\[4pt]
\mathcal{V} &\equiv 1/n, && \text{specific volume, i.e.\ volume per baryon, defined by}
\tag{2.5.3} \\[4pt]
\mathcal{V}_0 &= 1/\rho_0, && \text{specific volume, i.e.\ volume per unit rest mass, defined by}
\tag{2.5.4}
\end{align}

\begin{align}
\rho &= \rho_0(1 + \Pi), && \text{total density of mass-energy, as measured in LRF.}
\tag{2.5.5}\\[4pt]
\Pi &\quad && \text{specific internal energy, defined by}
\tag*{}
\end{align}

\textit{Note}: here and throughout this section we set the speed of light equal to
one.
\begin{align}
p  &\quad && \text{isotropic pressure, as measured in LRF.}
\tag*{} \\
T  &\quad && \text{temperature, as measured in LRF.}
\tag*{} \\
s  &\quad && \text{entropy per baryon, as measured in LRF; of course,}
\tag*{} \\
s_0 &= m_B^{-1}\,s. && \text{entropy per unit mass, as measured in LRF.}
\tag{2.5.6}
\end{align}

$\mu$ \quad chemical potential, as measured in LRF; defined by
%
\begin{equation}
\tag{2.5.7}
\mu \equiv \left(\frac{\partial\rho}{\partial n}\right)_{\!s}
  = \frac{\rho + p}{n}\,.
\end{equation}

(The second equality follows from the first law of thermodynamics,
below.)

\medskip
\noindent
$\mathbf{q}$ \quad flux of energy (due to heat conduction, radiation,
convection, etc.)\ as measured in LRF.  This flux
(ergs\,cm$^{-2}$\,sec$^{-1}$) is a purely spatial vector as measured
in LRF; i.e.
%
\begin{equation}
\tag{2.5.8}
\mathbf{q} \cdot \mathbf{u} = 0.
\end{equation}

\noindent
$\mathbf{S}$ \quad entropy density-flux vector.  In LRF this vector has
time component equal to the entropy density, $S^0 = ns$, and space
components equal to the entropy flux, $S^j = q^j / T$.  Hence, in
frame-independent notation
%
\begin{equation}
\tag{2.5.9}
\mathbf{S} = ns\,\mathbf{u} + \mathbf{q}/T.
\end{equation}

\noindent
Note that
%
\begin{equation}
\tag{2.5.10}
\boldsymbol{\nabla}\cdot\mathbf{S}
= \left\{
\begin{array}{l}
\text{rate at which entropy is being generated}\\
\text{per unit volume as measured in LRF}
\end{array}
\right\}.
\end{equation}

\bigskip
\noindent
\emph{First law of thermodynamics}

\smallskip
\noindent
Follow a fluid element, containing $A$ baryons, along its world tube.
The total mass-energy in the fluid element, $\rho A/n$, changes as a
result of compression (change in volume, $A/n$) and as a result of
influx of heat:
%
\begin{equation}
\tag{2.5.11}
d(\rho A/n) = -p\,d(A/n) + T\,d(As).
\end{equation}

\noindent
[Here and throughout this section we assume for simplicity no internal
generation of entropy in the fluid element---e.g., no irreversible
chemical reactions; this enables us to write the influx of heat in
terms of the change in entropy, $T\,d(As)$.]

One often uses $A = \text{const.}$ and $\mu \equiv (\rho + p)/n$ to
rewrite this first law of thermodynamics in the equivalent forms
%
\begin{align}
\tag{2.5.12}
d\rho &= \frac{\rho + p}{n}\, dn + nT\, ds, \\[4pt]
\tag{2.5.12$'$}
d\mu  &= V\, dp + T\, ds.
\end{align}

\noindent
From the first law one can read off partial derivatives, e.g.
%
\begin{equation}
\tag{2.5.12$''$}
(\partial\rho/\partial n)_s = (\rho + p)/n, \qquad
(\partial\mu/\partial p)_s = T.
\end{equation}

\bigskip
\noindent
\emph{Fundamental relation and equations of state}

\smallskip
\noindent
The peculiar thermodynamic properties of the particular fluid being
studied are determined by a ``fundamental thermodynamic relation''
%
\begin{equation}
\tag{2.5.13}
\rho = \rho(n,\,s) \qquad \text{or} \qquad
\mu = \mu(p,\,s).
\end{equation}

\noindent
Once this relation has been specified---and once the constant $m_B$ has
been specified---one can use the first law of thermodynamics
(2.5.12) or (2.5.12$'$) to derive explicit expressions
for all other thermodynamic variables as functions of $n$, $s$ or $p$,
$s$.  For example, by combining with the first law one can derive the
``equations of state''
%
\begin{equation}
\tag{2.5.8$'$}
T(n,s) = \frac{1}{n}\left(\frac{\partial\rho}{\partial s}\right)_{\!n},
\qquad
p(n,s) = n\!\left(\frac{\partial\rho}{\partial n}\right)_{\!s} - \rho.
\end{equation}

\bigskip
\noindent
\emph{Adiabatic index and speed of sound}

\smallskip
\noindent
One defines the adiabatic index $\Gamma_1$ by
%
\begin{equation}
\tag{2.5.14}
\Gamma_1
\equiv \left(\frac{\partial\ln p}{\partial\ln n}\right)_{\!s}
= \frac{\rho + p}{p}\,
  \left(\frac{\partial p}{\partial\rho}\right)_{\!s}\,,
\end{equation}
%
where the third equality follows from the first law of thermodynamics.
It turns out that weak adiabatic perturbations (weak ``sound waves'')
propagate, in the LRF, with ordinary velocity
%
\begin{equation}
\tag{2.5.15}
c_s
= \left[\left(\frac{\partial p}{\partial\rho}\right)_{\!s}\right]^{1/2}
= \left(\frac{\Gamma_1\, p}{\rho + p}\right)^{\!1/2}.
\end{equation}

\bigskip
\noindent
\emph{Second law of thermodynamics}

\smallskip
\noindent
Equation~(2.5.10) allows one to write the second law of thermodynamics
in the form
%
\begin{equation}
\tag{2.5.16}
\boldsymbol{\nabla}\cdot\mathbf{S} \geq 0.
\end{equation}

\bigskip
\noindent
\emph{Decomposition of 4-velocity}

\smallskip
\noindent
One decomposes the gradient of the 4-velocity, $\nabla\mathbf{u}$,
into its ``irreducible'' tensorial parts:
%
\begin{equation}
\tag{2.5.17}
u_{\alpha;\beta}
= \sigma_{\alpha\beta} + \omega_{\alpha\beta}
  + \tfrac{1}{3}\,\theta\, h_{\alpha\beta}
  - a_\alpha\, u_\beta\,.
\end{equation}

\noindent
Here $\mathbf{a}$ is the 4-acceleration of the fluid
%
\begin{equation}
\tag{2.5.18a}
\mathbf{a} \equiv \nabla_{\mathbf{u}}\,\mathbf{u},
\qquad\text{i.e.}\quad
a^\alpha = u^\alpha{}_{;\beta}\, u^\beta\,,
\end{equation}
%
and $-a_\alpha u_\beta$ is that portion of $u_{\alpha;\beta}$ which is
not orthogonal to $\mathbf{u}$.  The remainder of $u_{\alpha;\beta}$
(i.e.\ $\sigma_{\alpha\beta} + \omega_{\alpha\beta} +
\tfrac{1}{3}\theta\, h_{\alpha\beta}$) is orthogonal to $\mathbf{u}$;
i.e.\ in the LRF it has only spatial components.  This orthogonal part
is decomposed into an isotropic part,
%
\begin{equation}
\tag{2.5.18b}
\theta \equiv \nabla \cdot \mathbf{u}
= u^\alpha{}_{;\alpha}\,,
\end{equation}
%
expansion, $\tfrac{1}{3}\theta h_{\alpha\beta}$, where $\theta$ is the
``expansion''; and $h_{\alpha\beta}$ is the ``projection tensor''
%
\begin{equation}
\tag{2.5.18c}
h_{\alpha\beta} \equiv g_{\alpha\beta} + u_\alpha\, u_\beta\,;
\end{equation}
%
plus a symmetric, trace-free ``shear''
%
\begin{equation}
\tag{2.5.18d}
\sigma_{\alpha\beta}
= \tfrac{1}{2}(u_{\alpha;\mu}\, h^\mu{}_\beta
              + u_{\beta;\mu}\, h^\mu{}_\alpha)
  - \tfrac{1}{3}\theta\, h_{\alpha\beta}\,;
\end{equation}
%
plus an antisymmetric ``rotation'' or ``vorticity''
%
\begin{equation}
\tag{2.5.18e}
\omega_{\alpha\beta}
= \tfrac{1}{2}(u_{\alpha;\mu}\, h^\mu{}_\beta
              - u_{\beta;\mu}\, h^\mu{}_\alpha).
\end{equation}

An observer in the LRF sees all fluid elements in his neighborhood to
move with low (nonrelativistic) velocities.  Let him use standard
Newtonian methods to calculate or measure the expansion $\theta$; shear
$\sigma^{\hat{j}}_{\;\hat{k}}$; and rotation $\omega^{\hat{j}}_{\;\hat{k}}$
of the fluid, and let a relativist calculate $\theta$,
$\sigma^{\hat{j}}_{\;\hat{k}}$, and
$\omega^{\hat{j}}_{\;\hat{k}}$ in the LRF from the above equations.
The two calculations will give the same answers.
[See, e.g., Ellis (1971).]

\bigskip
\noindent
\emph{Frozen-in magnetic field}

\smallskip
\noindent
Consider an ionized plasma which contains a ``frozen-in'' magnetic field.
The field is pure magnetic (no electric field) in the LRF.  Therefore
it can be described by a magnetic-field 4-vector $\mathbf{B}$
orthogonal to $\mathbf{u}$,
%
\begin{equation}
\tag{2.5.19}
\mathbf{B}\cdot\mathbf{u} = 0.
\end{equation}

\noindent
[For simplicity we assume that the magnetic permeability of the plasma
is the same as that of vacuum; so we do not distinguish between
``$\mathbf{B}$'' and ``$\mathbf{M}$''.  For a more general treatment, see
Lichnerowicz (1967).]  As the fluid moves, carrying with it the
frozen-in B-field, $\mathbf{B}$ must change as
%
\begin{equation}
\tag{2.5.20}
\frac{DB^\alpha}{d\tau}
= u_{\alpha;\beta}\, B^\beta
  + (\sigma_{\alpha\beta} - \tfrac{2}{3}\theta\, h_{\alpha\beta})\, B^\beta\,.
\end{equation}

\noindent
The term $u_{\alpha;\beta}\, B^\beta$ is required to keep $\mathbf{B}$
orthogonal to $\mathbf{u}$; by the term $\sigma_{\alpha\beta}\, B^\beta$,
the rotation of the fluid rotates the field lines; by the term
$-\tfrac{2}{3}\theta\, h_{\alpha\beta}\, B^\beta$, the compression of the
fluid orthogonal to $\mathbf{B}$ magnifies $\mathbf{B}$, conserving
flux.

\bigskip
\noindent
\emph{Stress-energy tensor}

\smallskip
\noindent
Consider a fluid with isotropic pressure, with shear and bulk viscosity,
with energy flowing between fluid elements, and with a frozen-in
magnetic field.  The stress-energy tensor for such a fluid is
%
\begin{align}
\tag{2.5.21}
T^{\alpha\beta}
&= \rho u^\alpha u^\beta + (p - \zeta\theta)\, h^{\alpha\beta}
   - 2\eta\, \sigma^{\alpha\beta}
   + q^\alpha u^\beta + u^\alpha q^\beta + \pi^{\alpha\beta}
\notag \\
&\quad
   + \frac{1}{8\pi}(B^2 u^\alpha u^\beta + B^2 h^{\alpha\beta}
   - 2B^\alpha B^\beta).
\end{align}

\noindent
The term $\rho u^\alpha u^\beta$ is the total density of mass-energy
(excluding only that of the frozen-in B-field) as measured in the LRF.
The term $p\, h^{\alpha\beta}$ is the isotropic pressure that would be
measured in the LRF if the gas were not changing volume (if $\theta$
were zero).  The quantities $\zeta$ and $\eta$ are the
\emph{coefficients of bulk viscosity and of dynamic viscosity},
respectively.  The term $-\zeta\theta\, h^{\alpha\beta}$ is the
isotropic viscous stress which resists isotropic expansion
($\theta > 0$) or compression ($\theta < 0$) of the fluid.
The term $-2\eta\, \sigma^{\alpha\beta}$ is the viscous shear stress
which resists shearing motions.  The term
$q^\alpha u^\beta + u^\alpha q^\beta$ is the energy density and
stresses associated with the flowing energy $\mathbf{q}$, as measured
in the LRF, are here neglected by comparison with $\rho u^\alpha u^\beta$
and $p\, h^{\alpha\beta}$.

(\emph{Note:} the energy density and stresses associated with the
flowing energy $\mathbf{q}$, as measured in the LRF, are here
neglected by comparison with $\rho u^\alpha u^\beta$ and
$p\, h^{\alpha\beta}$.  This neglect is valid with enormous accuracy in
most contexts of interest in these lectures.  An exception is the
radiation pressure which produces ``self-regulation'' of accretion onto
black holes when the inflowing mass is sufficiently large; see
\S\S\,4.5 and 5.13.)  The term
%
\begin{equation}
\tag{2.5.22}
T^{\alpha\beta}_{\text{MAG}}
= \frac{1}{8\pi}(B^2 u^\alpha u^\beta + B^2 h^{\alpha\beta}
  - 2B^\alpha B^\beta)
\end{equation}
%
is the Maxwell stress-energy associated with the frozen-in B-field (in
LRF: energy density $B^2/8\pi$, pressure $B^2/8\pi$ orthogonal to
field lines; tension $-B^2/8\pi$ along field lines).

\bigskip
\noindent
\emph{Equations of hydrodynamics}

\smallskip
\noindent
The fundamental equations governing the motion of a fluid in a given
gravitational field (spacetime geometry) are (i) the law of
\emph{baryon conservation}
%
\begin{equation}
\tag{2.5.23}
\nabla \cdot (n\,\mathbf{u}) = 0,
\quad\text{i.e.}\quad dn/d\tau = -\theta\, n,
\end{equation}
%
or equivalently rest-mass conservation
%
\begin{equation}
\tag{2.5.23$'$}
\nabla \cdot (\rho_0\, \mathbf{u}) = 0,
\quad\text{i.e.}\quad d\rho_0/d\tau = -\theta\, \rho_0\,;
\end{equation}

\noindent
(ii) the law of \emph{local energy conservation}
%
\begin{equation}
\tag{2.5.24}
\mathbf{u}\cdot(\nabla\cdot\mathbf{T}) = 0;
\end{equation}

\noindent
(iii) the \emph{Euler equations} (i.e.\ law of local momentum
conservation)
%
\begin{equation}
\tag{2.5.25}
\mathbf{h}\cdot(\nabla\cdot\mathbf{T}) = 0;
\end{equation}

\noindent
(iv) the laws of thermodynamics, (2.5.11)--(2.5.16); (v) the laws of
energy transport for the frozen-in B-field, (2.5.20); the laws of
energy transport, which govern $\mathbf{q}$ (see \S\,2.6, 2.6.43).

\bigskip
\noindent
\emph{Law of local energy conservation}

\smallskip
\noindent
When one evaluates the law of local energy conservation,
$\mathbf{u}\cdot(\nabla\cdot\mathbf{T}) = 0$,
for the stress-energy tensor (2.5.21), one finds that the Maxwell
stress-energy
%
\begin{equation}
\tag{2.5.26}
\mathbf{u}\cdot(\nabla\cdot\mathbf{T}_{\text{MAG}}) = 0
\end{equation}
%
gives zero contribution (work done to compress magnetic field is
precisely equal to increase in magnetic field energy); and that the
remainder of the stress-energy tensor gives
%
\begin{equation}
\tag{2.5.27}
d\rho/d\tau = -(\rho + p)\,\theta
+ \zeta\,\theta^2 + 2\eta\,\sigma_{\alpha\beta}\,\sigma^{\alpha\beta}
- \nabla\cdot\mathbf{q} - \mathbf{a}\cdot\mathbf{q}\,.
\end{equation}

\noindent
Here $d/d\tau$ is derivative with respect to proper time along the
world lines of the fluid; i.e., in the language of the differential
geometer, $d/d\tau = \mathbf{u}$.  The terms $-(\rho + p)\,\theta$ is
the increase in mass-energy density due to compression.  The term
$\zeta\,\theta^2$ and $2\eta\,\sigma_{\alpha\beta}\,\sigma^{\alpha\beta}$
are the increases in mass-energy density due to viscous heating
(conversion of relative kinetic energy of adjacent fluid elements into
heat).  The term $-\nabla\cdot\mathbf{q}$ is the influx of
mass-energy from neighboring fluid elements.  The term
$-\mathbf{a}\cdot\mathbf{q}$ is a special relativistic correction to
the ``redshift'' of $\mathbf{q}$.  To understand this term, consider a
flux of energy $\mathbf{q}$---or, equivalently, with the inertia of
the energy flux $\mathbf{q}$ associated with it, a flux of mass-energy
from the viewpoint of an accelerated fluid (photons do not interact with
them).  The acceleration gives rise to a redshift (a gravitational
redshift) and hence to a nonzero $\nabla\cdot\mathbf{q}$ (no
interaction between fluid and photons), one must include the correction
factor $-\mathbf{a}\cdot\mathbf{q}$.  For a similar reason, a similar
relativistic correction appears in the law of heat conduction.

\bigskip
\noindent
\emph{Euler equation for a perfect fluid}

\smallskip
\noindent
Consider a perfect fluid, flowing adiabatically through spacetime
($\zeta = \eta = \mathbf{B} = \mathbf{q} = 0$).  In this case the
Euler equations $\mathbf{h}\cdot(\nabla\cdot\mathbf{T}) = 0$ reduce to
%
\begin{equation}
\tag{2.5.30}
(\rho + p)\,\mathbf{a} = -\nabla p - \mathbf{u}\,(\mathbf{u}\cdot\nabla p).
\end{equation}
%
In words:
\[
(\rho + p)\,\mathbf{a}
= -\left(\!\text{inertial mass per}\atop\text{unit volume}\!\right)
  \times\text{(4-acceleration)}
= -\left(\!\text{pressure gradient}\atop\text{projected orthogonal to }\mathbf{u}\!\right).
\]

\textit{Bernoulli equation}

Consider a perfect fluid undergoing stationary, adiabatic flow in a stationary
spacetime.  More particularly, set $\zeta = \eta = \fvec{B} = \fvec{q} = 0$; assume that spacetime is
endowed with a Killing vector field $\boldsymbol{\xi}$,

\begin{equation}
\xi_{\alpha;\beta} + \xi_{\beta;\alpha} = 0\,,
\tag{2.5.31}
\end{equation}

which need not be timelike; and assume that the flow is adiabatic, $ds/d\tau = 0$,
and stationary in the sense that

\begin{equation}
\mathscr{L}_{\boldsymbol{\xi}}\,p = 0, \quad
\nabla_{\boldsymbol{\xi}}\,\rho = \nabla_{\boldsymbol{\xi}}\,p = \cdots = 0\,.
\tag{2.5.32}
\end{equation}

Here $\mathscr{L}_{\boldsymbol{\xi}}$ is the Lie derivative along $\boldsymbol{\xi}$.  In this case the Euler equations~(2.5.30),
together with the first law of thermodynamics~(2.5.12), imply the relativistic
Bernoulli equation

\begin{equation}
d(\mu\,\fvec{u}\cdot\boldsymbol{\xi})/d\tau = 0\,;
\tag{2.5.33}
\end{equation}

i.e., $\mu\,\fvec{u}\cdot\boldsymbol{\xi}$ is constant along flow lines.

\bigskip
\textit{Newtonian limit}

The Newtonian limit of relativistic hydrodynamics is obtained when, in a nearly
global Lorentz frame, the following approximations hold:

\begin{gather}
g_{00} = -(1 + 2\Phi), \quad |\Phi| \ll 1, \quad
B^2/\rho_0 \ll 1,
\notag\\
p/\rho_0 \ll 1, \quad \Pi \ll 1, \quad v^2 \ll 1\,.
\tag{2.5.34}
\end{gather}

Here $\Phi$ is the Newtonian gravitational potential with sign $\Phi < 0$, and $v$ is the
ordinary velocity of the fluid

\begin{equation}
\fvec{v} \equiv v_j = (u^j/u^0)\,\fvec{e}_j\,.
\tag{2.5.35}
\end{equation}

In the Newtonian limit the chemical potential $\mu$ reduces to

\begin{equation}
\mu = m_B(1 + w)\,,
\tag{2.5.36}
\end{equation}

%%% ============================================================
%%% Book page 370
%%% (PDF p.15, left column)
%%% ============================================================

where $w$ is the \textit{enthalpy}

\begin{equation}
w = \Pi + p/\rho_0\,.
\tag{2.5.37}
\end{equation}

The Bernoulli equation~(2.5.33) reduces to the familiar form

\begin{equation}
\Phi + \tfrac{1}{2}v^2 + w = \text{constant along flow lines}\,.
\tag{2.5.38}
\end{equation}

The Euler equation for a perfect fluid in adiabatic flow, (2.5.30), reduces to

\begin{equation}
\frac{d\fvec{v}}{dt} = -\nabla\Phi - \frac{1}{\rho_0}\nabla p\,,
\tag{2.5.39}
\end{equation}

where $d/dt$, the derivative with respect to proper time along the flow lines, has
the Newtonian form

\begin{equation}
d/dt = \partial/\partial t + \fvec{v}\cdot\boldsymbol{\nabla}\,.
\tag{2.5.40}
\end{equation}

The first law of thermodynamics, (2.5.12) and (2.5.12$'$), reduces to

\begin{equation}
d\Pi = -p\,dV_0 + T\,ds_0\,,
\tag{2.5.41}
\end{equation}

\begin{equation}
dw = V_0\,dp + T\,ds_0\,.
\tag{2.5.42}
\end{equation}

(In Newtonian theory one usually adopts a per-unit-mass viewpoint rather than a
per-baryon viewpoint; and thus one uses $\rho_0$, $V_0$ rather than $n$, $\mathcal{V}$, $s_0$.)

%%% ============================================================
%%% Section 2.6: Radiative Transfer
%%% Book page 370–371 (PDF p.15, left–right columns)
%%% ============================================================

\subsection{Radiative Transfer}
\label{sec:2.6}

\textit{Basic references}\quad
Chandrasekhar (1960); Appendix~1 of Pacholczyk (1970); Mihalas (1970).

\bigskip
\textit{Notation and terminology}

At an arbitrary event in spacetime pick an arbitrary
local Lorentz frame.  In that frame pick an arbitrary spatial direction $\fvec{n}$ (a
unit vector; see Figure~2.6.1).  Examine the amount of energy $dE$ that is (i) carried
by photons across a unit surface area $dA$ orthogonal to $\fvec{n}$ (the surface $\mathscr{S}_m$ of
Figure~2.6.1) during unit time $dt$, with (ii) the photons having frequencies $\nu$ in
the range $d\nu$, and (iii) the photons being directed into a solid angle $d\Omega$ about
$\fvec{n}$.  The ratio

\begin{equation}
I_\nu = \frac{dE}{dt\;dA\;d\nu\;d\Omega}
\tag{2.6.1}
\end{equation}

is called the \textit{specific intensity} or simply the \textit{intensity}.  It depends on (a) the
event in spacetime; (b) the choice of Lorentz frame; (c) the direction $\fvec{n}$; (d)
the frequency $\nu$.  One also defines the \textit{total intensity} $I$ by

\begin{equation}
I = \int I_\nu\;d\nu = \frac{dE}{dt\;dA\;d\Omega}\,;
\tag{2.6.2}
\end{equation}

the \textit{average specific intensity} or \textit{average intensity} $J_\nu$ (averaged over all directions)
by

\begin{equation}
J_\nu = \frac{1}{4\pi}\int I_\nu\;d\Omega\,;
\tag{2.6.3}
\end{equation}

and the \textit{average total intensity} $J$ by

\begin{equation}
J = \frac{1}{4\pi}\int I\;d\Omega = \frac{1}{4\pi}\int I_\nu\;d\nu\;d\Omega\,.
\tag{2.6.4}
\end{equation}

It is easy to see that the \textit{energy density per unit frequency} in the radiation is

\begin{equation}
\rho^{(\text{rad})} = \frac{dE}{d^3x\;d\nu} = \frac{4\pi J_\nu}{c}\,,
\tag{2.6.5}
\end{equation}

and that the \textit{total energy density} in the radiation is

\begin{equation}
\rho^{(\text{rad})} = \frac{dE}{d^3x} = \frac{4\pi J}{c}\,.
\tag{2.6.6}
\end{equation}

%%% ============================================================
%%% Book page 371 (right column) — Figure 2.6.1, flux, invariance
%%% ============================================================

Pick a 2-surface $\mathscr{S}_m$ in the chosen Lorentz frame, with unit normal $\fvec{m}$ (Fig.\
2.6.1).  Examine the total energy $dE$ per unit area $dA$ that is carried across the
chosen surface during unit time $dt$, by photons having frequency $\nu$ in the range
$d\nu$.  Make no restrictions on the photon direction or solid angle; but count
negatively those photons that cross $dA$ from the front side toward the back.
The ratio

\begin{equation}
F_\nu = \frac{dE}{dt\;dA\;d\nu}
\tag{2.6.7}
\end{equation}

is called the \textit{specific flux}, or simply the \textit{flux}.  It is easy to see from Figure~2.6.1
that

\begin{equation}
F_\nu = \int I_\nu\,\cos\theta\;d\Omega\,.
\tag{2.6.8}
\end{equation}

\begin{figure}[htbp]
\centering
\fbox{\parbox{0.9\columnwidth}{\centering [Figure 2.6.1 placeholder]\\
The surfaces, directions, and solid angle used in the definition of intensity
and of flux.  $\mathscr{S}_m$ is a surface with unit normal $\fvec{m}$;
$\fvec{n}$ is a unit vector in the photon propagation direction;
$\theta$ is the angle between $\fvec{m}$ and $\fvec{n}$; $d\Omega$ is the solid angle element.}}
\caption{The surfaces, directions, and solid angle used in the definition of intensity
and of flux.}
\label{fig:2.6.1}
\end{figure}

The specific flux depends on (a) the chosen event in spacetime; (b) the chosen
Lorentz frame at that event; (c) the chosen 2-surface $\mathscr{S}_m$ in that Lorentz frame.

%%% ============================================================
%%% Book page 372 — Invariance of I_nu/nu^3, equation of radiative transfer
%%% (PDF p.16, left column)
%%% ============================================================

The integral of the flux over all frequencies is called the \textit{total flux}

\begin{equation}
F = \int F_\nu\;d\nu = \int I\,\cos\theta\;d\Omega\,.
\tag{2.6.9}
\end{equation}

Notice that the total flux is the ``first moment'' of the total intensity (the
``moment'' being taken with respect to the unit normal $\fvec{m}$).  One can readily
verify that the ``second moment'' is equal to the radiation pressure that acts
across the surface $\mathscr{S}_m$:

\begin{equation}
T^{(\text{rad})} \equiv \fvec{m}\cdot\fvec{T}^{(\text{rad})}\cdot\fvec{m}
= \int I\,\cos^2\theta\;d\Omega\,.
\tag{2.6.10}
\end{equation}

\bigskip
\textit{Invariance of $I_\nu/\nu^3$}

Choose a particular photon that passes through a given event in spacetime.
Observe that particular photon, and all others in the vicinity, from several different
Lorentz frames at the given event.  The frequency $\nu$ of the photon will depend
on the choice of Lorentz frame (Doppler shift from one frame to another).
The specific intensity $I_\nu$ in the frame neighborhood of the photon will also depend
on the choice of Lorentz frame.  However, the ratio $I_\nu/\nu^3$ will be Lorentz-
invariant. (Aside from a factor $h^{-4}$, $I_\nu/\nu^3$ is the invariant number density in
phase space

\begin{equation}
\mathscr{N} = \frac{dN}{d^3x\,d^3p} = \frac{1}{h^3}\,\frac{I_\nu}{\nu^3}\,;
\tag{2.6.11}
\end{equation}

see, e.g., \S22.6 of MTW.)

When photons propagate freely through curved spacetime (no interaction with
matter), the ratio $I_\nu/\nu^3$ is conserved along the world line of each photon
(Liouville's theorem).

\paragraph*{Equation of Radiative Transfer}

Consider the propagation of photons through a medium (e.g., through gas that
is falling into a black hole).\ Analyze the photon propagation from the viewpoint
of observers at rest in the medium (local rest frame; ``LRF'').\ At each event
denote by $\mathbf{u}$ the 4-velocity of the medium---and hence also of the LRF.\ Focus
attention on all photons in the (phase-space) neighborhood of a given null
geodesic ray.\ Denote by $\mathbf{p}$ the 4-momentum of that ray.\ Then at each event
along the given ray $\mathbf{p}$ is given by
%
\begin{equation}
\mathbf{p} = h\nu(\mathbf{u} + \mathbf{n}).
\tag{2.6.12}
\end{equation}
%
Here $\mathbf{n}$ is a unit vector;
%
\[
\mathbf{n} \cdot \mathbf{u} = 0,
\]
%
and (i) is interpreted in the LRF as the direction of propagation of the ray
(same as 3-vector $\mathbf{n}$ of Figure~2.6.1).\ Also, in eq.~(2.6.12), $\nu$ is the frequency of
any photon which propagates along the ray, and $h\nu$ is its energy, as measured in
the LRF.

Because of the acceleration and shear of the medium, the frequency
$\nu$ of the chosen ray changes from point to point along the ray.\ The change $d\nu$,
when the ray propagates a proper spatial distance $dl$ as seen in the LRF, is
%
\begin{equation}
d\nu = \nu(-\,\mathbf{p} \cdot \mathbf{u}/h)\,dl.
\tag{2.6.13}
\end{equation}
%
(This frame-independent equation is derived easily by geometrical arguments
in the LRF.)\ A straightforward calculation, using the geodesic equation for the
ray
%
\[
\boldsymbol{\nabla}_{\mathbf{p}}\,\mathbf{p} = 0,
\]
%
and using expansion (2.5.17) for $\boldsymbol{\nabla}\mathbf{u}$, reveals
%
\begin{equation}
d\nu/dl = -\nu(\mathbf{n} \cdot \mathbf{a} + \tfrac{1}{3}\theta
  + n^\alpha n^\beta \sigma_{\alpha\beta}).
\tag{2.6.14}
\end{equation}
%
The first term, $-\nu\,\mathbf{n} \cdot \mathbf{a}$, is the ``gravitational redshift'' produced by the acceleration
of the LRF; the second and third terms, $-\nu(\tfrac{1}{3}\theta + n^\alpha n^\beta \sigma_{\alpha\beta})$, are the
``cosmological redshift'' due to the expansion of the medium along the direction
of the ray.

If there were no interaction with the medium, then $I_\nu/\nu^3$ would be conserved
along the ray.\ Hence, the equation of radiative transfer along the given ray must
have the form
%
\begin{equation}
\frac{dI_\nu}{dl} - \left(\frac{3}{\nu}\,\frac{d\nu}{dl}\right) I_\nu
  = \text{(effects of interaction with medium)};
\tag{2.6.15}
\end{equation}
%
i.e.,
%
\[
dI_\nu/dl + (3\mathbf{n} \cdot \mathbf{a} + \theta
  + 3n^\alpha n^\beta \sigma_{\alpha\beta})\,I_\nu
  = \text{(interaction effects)}.
\]
%
Four types of interaction can occur: spontaneous emission of radiation by the
matter; stimulated emission; absorption; and scattering.

We shall assume that on all length scales of interest the medium is isotropic;
and we shall denote its emissivity by
%
\begin{equation}
\varepsilon_\nu \equiv \frac{dE}{d\nu\,dt\,dm_0}
  = \begin{pmatrix} \text{energy emitted spontaneously per unit frequency}\\
    \text{during unit time by a unit rest mass, integrated}\\
    \text{over all angles, and measured in LRF} \end{pmatrix}.
\tag{2.6.16}
\end{equation}
%
(The emissivities for free-free and free-bound transitions were discussed in \S\S\,2.1
and 2.2.)\ Then spontaneous emission contributes
%
\begin{equation}
\left(\frac{dI_\nu}{dl}\right)_{\!\text{spontaneous emission}}
  = \frac{1}{4\pi}\,\rho_0\,\varepsilon_\nu
\tag{2.6.17}
\end{equation}
%
to the specific intensity along the given ray.\ Here $\rho_0$ is the density of rest mass
in the medium.

The rate of absorption and the rate of stimulated emission are both proportional
to the intensity of the passing beam.\ Hence, the effects of absorption and of
stimulated emission can be lumped together into a single \textit{absorption coefficient}
$\kappa_\nu$, defined by
%
\begin{equation}
\left(\frac{dI_\nu}{dl}\right)_{\!\substack{\text{absorption plus}\\[2pt]
  \text{stimulated emission}}}
  = -\rho_0\,\kappa_\nu\,I_\nu.
\tag{2.6.18}
\end{equation}
%
The dimensions of the absorption coefficient are $\mathrm{cm}^2/\mathrm{g}$ (``absorption cross section
per unit mass'' at the given frequency).\ When absorption dominates over
stimulated emission, $\kappa_\nu$ is positive; when stimulated emission dominates,
$\kappa_\nu$ is negative (``negative absorption'').

Scattering is more complicated to treat than emission and absorption.\ Let
%
\begin{equation}
\frac{d\kappa_s}{d\Omega'\,d\nu'}\,(\mathbf{p},\,\mathbf{p}')
\tag{2.6.19}
\end{equation}
%
be the differential cross section (as measured in the LRF) for a unit rest mass to
convert a photon of momentum $\mathbf{p}$ into a photon of momentum $\mathbf{p}'$; and let
%
\[
\kappa_s(\nu) = \int \frac{d\kappa_s}{d\Omega'\,d\nu'}\,d\Omega'\,d\nu'
\]
%
be the total scattering cross-section (``scattering opacity'') per unit rest mass.
For example, in the case of electron scattering by a nonrelativistic plasma
($h\nu \ll m_e c^2$; $kT \ll m_e c^2$; see \S\,4.4), when one neglects the tiny change in photon
frequency the cross sections per unit rest mass are
%
\begin{equation}
\frac{d\kappa_s}{d\Omega'\,d\nu'}\,(\mathbf{p},\,\mathbf{p}')
  = \frac{f_e}{2m_p}\,\sigma_T\,\tfrac{3}{8\pi}\,
    [1 + (\mathbf{n} \cdot \mathbf{n}')^2]\,\delta(\nu' - \nu)
\tag{2.6.20}
\end{equation}
%
\[
\kappa_s = (f_e/m_p)\,\sigma_T = (0.40\;\mathrm{cm}^2/\mathrm{g})\,f_e
\]
%
(Recall: $n = \rho/h\nu$, $n' = \rho'/h\nu'$; $f_e$ is the number of free electrons per baryon in
the plasma.)\ Scattering, as described by the appropriate differential cross section,
can increase $I_\nu$ (scattering into the beam from other directions), or can decrease
$I_\nu$ (scattering out of the beam into other directions):
%
\begin{equation}
\left(\frac{dI_\nu}{dl}\right)_{\!\text{scattering}}
  = +\rho_0 \!\int \frac{d\kappa_s}{d\Omega\,d\nu}\,
    (\mathbf{p}',\,\mathbf{p})\,I_{\nu'}\,d\Omega'\,d\nu'
    \;-\; \rho_0\,\kappa_s\,I_\nu.
\tag{2.6.21}
\end{equation}
%
By combining equations (2.6.15)--(2.6.21), we obtain our final form of the
equation of radiative transfer
%
\begin{align}
dI_\nu/dl &+ (3\mathbf{n} \cdot \mathbf{a} + \theta
  + 3n^\alpha n^\beta \sigma_{\alpha\beta})\,I_\nu \notag\\
  &= +\rho_0 \!\int \frac{d\kappa_s}{d\Omega\,d\nu}\,
     (\mathbf{p}',\,\mathbf{p})\,I_{\nu'}\,d\Omega'\,d\nu'
     - \rho_0\,\kappa_s\,I_\nu \notag\\
  &\quad + \tfrac{1}{4\pi}\,\rho_0\,\varepsilon_\nu
     - \rho_0\,\kappa_\nu\,I_\nu.
\tag{2.6.22}
\end{align}

%%% -----------------------------------------------------------------
%%% Relationship between emissivity and absorption
%%% -----------------------------------------------------------------

\paragraph*{Relationship between emissivity and absorption}

Into an insulated box place a material medium and a radiation field; and then wait
until complete thermodynamic equilibrium is achieved.\ If $T$ is the equilibrium
temperature, then the radiation field as seen in the LRF will be isotropic with the
standard black-body intensity, $I_\nu = B_\nu$, where
%
\begin{equation}
B_\nu = \frac{2h\nu^3/c^2}{e^{h\nu/kT} - 1}.
\tag{2.6.23}
\end{equation}
%
In equilibrium there must be no net change of $I_\nu$ along any ray: scattering into
the beam must be completely balanced by scattering out of the beam; and
emission must be completely balanced by absorption.\ The requirement of
``detailed balance'' for emission and absorption can be met only if the emissivity
and the absorption coefficient are related by
%
\begin{equation}
\frac{dI_\nu}{dl} = \rho_0\,\kappa_\nu
  \!\left(\frac{\varepsilon_\nu}{4\pi\kappa_\nu}
  - I_\nu\right) = 0;
\tag{2.6.24}
\end{equation}
%
i.e.,
%
\[
\varepsilon_\nu / 4\pi\kappa_\nu = B_\nu.
\]
%
Since $\varepsilon_\nu$ and $\kappa_\nu$ depend only on the thermodynamic state of the matter, and
have nothing to do with the state of the radiation, relation~(2.6.24) must be
satisfied not only inside our insulated box, but also in all other cases where
the matter by itself is in thermodynamic equilibrium but the radiation might
not be.

For matter not in thermodynamic equilibrium one can use a more sophisticated
version of the principle of detailed balance (``Einstein A and B coefficients'') to
derive a more complicated relationship between the emissivity and the absorption
coefficient.\ See, e.g., Chandrasekhar~\cite{Chandrasekhar1960}.

As an application of the equilibrium relationship $\varepsilon_\nu/4\pi\kappa_\nu = B_\nu$, consider free-free
transitions in an ionized gas.\ Whenever the free electrons (which do the
emitting and absorbing) are in thermodynamic equilibrium with each other
(Maxwell--Boltzmann velocity distribution), the absorption coefficient---as
derived from the emissivity (2.1.27)---must be
%
\begin{equation}
\kappa_\nu^{\mathrm{ff}} = (1.50 \times 10^{25}\;\mathrm{cm}^5\;\mathrm{g}^{-1})
  \left(\frac{\rho_0}{\mathcal{A}}\right)
  \left(\frac{(eZ)^2}{\mathcal{A}}\right)
  \bar{G}\;T^{-7/2}\;\nu^{-3}
  \left(\frac{1 - e^{-x}}{x^{-1}}\right),
\tag{2.6.25}
\end{equation}
%
where
\[
x = h\nu/kT.
\]
%
See \S\,2.1 for notation.

Notice that for matter in thermodynamic equilibrium the relationship
$\varepsilon_\nu / 4\pi\kappa_\nu = B_\nu$ permits one to rewrite the equation
of radiative transfer~(2.6.22) in the form
%
\begin{equation}
\tag{2.6.26}
\frac{dI_\nu}{dl} + (3n\cdot a + \dot\theta + 3\sigma^{\alpha\beta}
n_\alpha n_\beta)\, I_\nu
= \rho_0\,\kappa_\nu\bigl(B_\nu - I_\nu\bigr)
  + \left.\frac{dI_\nu}{dl}\right|_{\text{scattering}}
  + \left.\frac{d(I_\nu/\nu^3)}{d\tau_\nu}\right|_{\text{scattering}} ;
\end{equation}
%
or, equivalently,
%
\begin{equation}
\tag{2.6.27}
\frac{d(I_\nu/\nu^3)}{d\tau_\nu}
= \rho_0\,\kappa_\nu \left(\frac{B_\nu}{\nu^3}
  - \frac{I_\nu}{\nu^3}\right)
  + \left.\frac{d(I_\nu/\nu^3)}{d\tau_\nu}\right|_{\text{scattering}}.
\end{equation}

%%% ---------------------------------------------------------------
%%% Optical depth
%%% ---------------------------------------------------------------

\bigskip
\noindent
\emph{Optical depth}

\medskip
\noindent
Consider radiation propagating out of a medium into surrounding empty space.
Follow a given ray ``backward'', from ``infinity'' into the medium.  Let the ray
have frequency $\nu_\infty$ at infinity; then its frequency at location~$l$
($l = {}$proper distance measured in LRF of medium) is
%
\begin{equation}
\tag{2.6.28}
\nu(l) = \nu_\infty \exp\!\left[\,-\int_i^l
\bigl(\dot\theta + n^\alpha n^\beta \sigma_{\alpha\beta}\bigr)\,dl\right].
\end{equation}
%
[cf.\ eq.~(2.6.14)].  In calculating the change of intensity along the ray, it is
often useful to replace the proper-length parameter~$l$ by the
``optical-depth'' parameter
%
\begin{equation}
\tag{2.6.29}
\tau_\nu = \int_i^l \rho_0\,\kappa_\nu\,dl.
\end{equation}

In the integration $\kappa_\nu$ must be evaluated at the frequency
$\nu(l)$.  The optical depth $\tau_\nu$ depends on (i) the world line of
the chosen ray, (ii) the frequency of the ray in spacetime, (iii) location
along that ray.  Then the equation of transfer
world line; and (iii) location along the ray at ``infinity'' $\nu_\infty$\ldots

One can also introduce an optical depth for scattering radiation out of the
beam, $\tau_s$:
%
\begin{equation}
\tag{2.6.30}
\tau_s = \int_i^l \rho_0\,\kappa_s\,dl.
\end{equation}

In terms of optical depths, the law of radiative transfer~(2.6.27) reads
%
\begin{equation}
\tag{2.6.31}
\frac{d(I_\nu/\nu^3)}{d\tau_\nu}
= \frac{B_\nu}{\nu^3} - \frac{I_\nu}{\nu^3}
  - \frac{\kappa_s}{\kappa_\nu}
  \left.\frac{d(I_\nu/\nu^3)}{d\tau_\nu}\right|_{\text{scattering}},
\end{equation}
%
where
%
\begin{equation}
\tag{2.6.32}
\left.\frac{d(I_\nu/\nu^3)}{d\tau_\nu}\right|_{\text{scattering}}
= \frac{\kappa_s}{\kappa_\nu}\left[\frac{I_\nu}{\nu^3} - \frac{1}{\kappa_s}
  \int \frac{dk_s}{d\Omega\,d\nu'}\,
  (\mathbf{p}',\mathbf{p})\,I_\nu'\,d\Omega'\,d\nu'
  - \rho_0\,\kappa_s\,I_\nu \right].
\end{equation}

%%% ---------------------------------------------------------------
%%% Optically thin / optically thick
%%% ---------------------------------------------------------------

One says that the medium is \emph{optically thin} if everywhere along
the ray $\tau_\nu \ll 1$ (or $\tau_s \ll 1$).

Consider a medium that (i) has its matter in thermodynamic equilibrium with
a spatially uniform temperature, (ii) is optically thick along a chosen ray, and
(iii) has absorption dominant over scattering, i.e.,
%
\begin{equation}
\tag{2.6.33}
\kappa_\nu \gg \kappa_s.
\end{equation}

One says that a medium is \emph{optically thick} to emission and absorption (or to
scattering) along a given ray if optical depths $\tau_\nu \gg 1$
(or $\tau_s \gg 1$) are achieved.

\medskip
\noindent
Along that ray, then the radiation emerging to infinity along the chosen ray must
have the blackbody form
%
\begin{equation}
\tag{2.6.34}
\frac{I_\nu}{\nu^3} = \frac{B_\nu}{\nu^3}
= \frac{2h\nu^3/c^3}{e^{h\nu/kT} - 1}.
\end{equation}
%
[One can prove this easily from the equation of transfer~(2.6.31).]\
Moreover, at a point in the medium where all rays are optically thick, the
radiation will have a blackbody intensity at all frequencies and in all
directions; so the energy density and pressure in the radiation will be
%
\begin{equation}
\tag{2.6.35}
\rho^{\text{rad}} = 3p^{\text{rad}} = bT^4,
\end{equation}
%
where $b$ is the universal constant
%
\begin{equation}
\tag{2.6.36}
b = \frac{8\pi^5 k^4}{15c^3 h^3}
  = 7.56 \times 10^{-15}\;\frac{\text{ergs}}{\text{cm}^3\,\text{K}^4}.
\end{equation}

Consider, alternatively, a medium that is optically thin along a chosen ray,
and has negligible scattering ($\tau_s \ll 1$) along that ray.  Then the
equation of transfer~(2.6.31) predicts for the radiation emerging to infinity
%
\begin{equation}
\tag{2.6.37}
I_\nu\,\nu^3 = \int_{\text{entire ray}} (\varepsilon_\nu/\nu^3)\,\rho_0\,dl
= (1/4\pi)\int (\varepsilon_\nu/\nu^3)_{\text{in region of strongest emission}}
  \rho_0\,dl.
\end{equation}

\begin{equation}
\tag{2.6.38}
\simeq \left(\int \rho_0\,dl\right)\!(1/4\pi)\,
(\varepsilon_\nu/\nu^3)_{\text{in region of strongest emission}}.
\end{equation}

In summary, the radiation from an optically thin source with negligible
scattering $I_\nu \propto \varepsilon_\nu$, and has an intensity proportional
to the amount of matter along the line of sight.  But radiation from an
optically thick source in thermodynamic equilibrium with negligible scattering
has the blackbody form independent of the nature of its emissivity.  Media with
non-negligible scattering will be studied in~\S5.10.

%%% ---------------------------------------------------------------
%%% Radiative transfer in the diffusion approximation
%%% ---------------------------------------------------------------

\bigskip
\noindent
\emph{Radiative transfer in the diffusion approximation}

\medskip
\noindent
Consider the interior of an optically thick medium $\tau_\nu \gg 1$ which is
in local thermodynamic equilibrium.  Suppose that the medium has a temperature
gradient and/or an acceleration, but that the characteristic length scale
%
\begin{equation}
\tag{2.6.39}
l_T \equiv \frac{T}{|\boldsymbol{\nabla}T + \dot{\mathbf{a}}\,T|}
\end{equation}
%
over which the temperature changes are long compared to the mean-free path of a
photon,
%
\begin{equation}
\tag{2.6.40}
l_T \gg 1/\kappa_\nu\,\rho_0 \equiv l_\nu.
\end{equation}
%
Then the radiation distribution will consist of a large, isotropic, blackbody
component, plus a tiny ``correction'' due to the temperature gradient
%
\begin{equation}
\tag{2.6.41}
I_\nu \simeq B_\nu + I_\nu^{(1)};\qquad
I_\nu^{(1)} \ll B_\nu.
\end{equation}

There is no net flux $F$ across any surface associated with the blackbody
component~$B_\nu$; but the ``correction'' term $I_\nu^{(1)}$ will lead to a
flux in the direction of
$\mathbf{h}\cdot(\boldsymbol{\nabla}T+\dot{\mathbf{a}}\,T)$.

(Here $\mathbf{h}$ is the projection operator of~\S2.5.)  One can calculate
the magnitude of that flux by inserting expression~(2.6.41) for~$I_\nu$ into
the equation of transfer~(2.6.22), and by then integrating over
$\cos\theta\,d\Omega\,d\nu$.  (Here $\theta$ is the angle between
$d\Omega$~and the direction of
$\mathbf{h}\cdot(\boldsymbol{\nabla}T+\dot{\mathbf{a}}\,T)$.)
The result is
%
\begin{equation}
\tag{2.6.43}
\mathbf{q} = (1/\kappa_R\rho_0)(\tfrac{1}{3}\partial T^4/\partial \mathbf{x})
  \,\mathbf{h}\cdot(\boldsymbol{\nabla}T + \dot{\mathbf{a}}\,T).
\end{equation}

Here $\mathbf{q}$ is the energy flux vector of~\S2.5, and its magnitude is the
flux of energy in the
$\mathbf{h}\cdot(\boldsymbol{\nabla}T+\dot{\mathbf{a}}\,T)$ direction
%
\begin{equation}
\tag{2.6.44}
|\mathbf{q}| = F.
\end{equation}

Also, in eq.~(2.6.43) $\bar{\kappa}$ is the ``Rosseland mean opacity'',
defined by
%
\begin{equation}
\tag{2.6.45}
\frac{1}{\bar{\kappa}}
\equiv \frac{\displaystyle\int_0^\infty
  (\kappa_\nu + \kappa_s)^{-1}\,
  (dB_\nu/dT)\,d\nu}
  {\displaystyle\int_0^\infty (dB_\nu/dT)\,d\nu}.
\end{equation}

For free-free transitions by themselves (eq.~2.6.25) the Rosseland mean
opacity is
%
\begin{equation}
\tag{2.6.46}
\kappa_{ff} = (0.645\times 10^{23}\;\text{cm}^2/\text{g})
  \left(\frac{C_f\,\bar{f}\,\bar{g}^2}{\mathcal{A}}\right)
  \left(\frac{\rho_0}{\text{g\;cm}^{-3}}\right) T^{-7/2}.
\end{equation}

See~\S2.1 for notation.  For electron scattering by itself in the nonrelativistic
case, the Rosseland mean opacity is
%
\begin{equation}
\tag{2.6.47}
\bar{\kappa}_{es} = \kappa_{es} = 0.40\;\text{cm}^2/\text{g}.
\end{equation}


%%% ===============================================================
%%% SECTION 2.7: Shock Waves
%%% ===============================================================

\subsection{Shock Waves}
\label{sec:2.7}

\emph{Basic references}\quad
Zel'dovich and Raizer~(1966)~\cite{ZeldovichRaizer1966};
Landau and Lifshitz~(1959)~\cite{LandauLifshitz1959};
Taub~(1948)~\cite{Taub1948};
Lichnerowicz~(1967, 1970, 1971)~\cite{Lichnerowicz1967};
Thorne~(1973a)~\cite{Thorne1973a}.

\bigskip
\noindent
\emph{Types of shock waves}

\medskip
\noindent
When gas in supersonic flow encounters an obstacle (e.g., the surface of a star,
or the geometric structure of the ergosphere of a black hole), a shock front
develops.  In the shock front the flow decelerates sharply from supersonic to
subsonic, and some of the kinetic energy of the flow gets converted into heat
(increase in entropy).

The structure of the shock depends on the nature of the forces which decelerate
the gas particles (atoms, ions, electrons).  If those forces are collisions between
the particles themselves (the usual case), one has an ordinary shock.  But if the
particles are decelerated without colliding---e.g., by impact onto the dipole
magnetic field of a neutron star, which swings the particles into Larmour
orbits---one has a \emph{collisionless} shock.  These notes will be confined to ordinary
shocks.  For the theory of collisionless shocks see, e.g., the end of~\S12 of Kaplan
(1966).

Shock waves in a partially ionized gas (e.g., the interstellar medium) can have
a somewhat different form than ordinary shocks.  As gas passes through the shock
front, much of its kinetic energy of supersonic flow can be converted into
excitation energy of atoms and ions, and can be quickly radiated away as ``line''
radiation.  The result is a bright glow from the shock front itself.  See, e.g.,
Pikel'ner~(1961) and~\S10 of Kaplan~(1966).  In this section we shall ignore such
energy losses to radiation---i.e., we shall demand that the gas behind the shock
have the same total energy per unit rest mass as the gas in front of the shock.

%%% ---------------------------------------------------------------
%%% The relativistic Rankine--Hugoniot equations
%%% ---------------------------------------------------------------

\bigskip
\noindent
\emph{The relativistic Rankine--Hugoniot equations}

\medskip
\noindent
Pick a particular event~$\mathscr{P}$ on a shock front.  In the neighborhood
of~$\mathscr{P}$ introduce a local Lorentz frame (``rest frame of the shock'')
in which (i) the shock is the surface $y = z = 0$; and (iii) on both sides of
the shock the fluid is moving in the $x$~direction, i.e., perpendicular to the
shock front (``normal shock''; see Figure~\ref{fig:2.7.1}).  That such a local
Lorentz frame exists in general one can prove quite easily [see, e.g.,
Taub~(1948)].

Denote the ``front'' side of the shock (side \emph{from} which the fluid moves)
by a ``1'', and denote the ``back'' side (side \emph{toward} which the fluid
moves) by a ``2''; see Figure~\ref{fig:2.7.1}.  Denote the velocity of the
fluid, as measured in the rest frame of the shock, by
%
\begin{subequations}
\begin{align}
\tag{2.7.1a}
v_1 &= (dx/dt)_1 = \text{ordinary velocity on front side},\\[4pt]
\tag{2.7.1b}
v_2 &= (dx/dt)_2 = \text{ordinary velocity on back side}.
\end{align}
\end{subequations}
%
(Note that these ``4-velocities'' are scalars, not vectors.)

In the rest frame of the shock the law of baryon conservation is equivalent
to continuity of the \emph{baryon flux}:
%
\begin{equation}
\tag{2.7.2a}
j \equiv n_1\,u_1 = n_2\,u_2,
\end{equation}
%
and the relations $n = 1/V_1$, $\mu_1 = p_1$,
%
\begin{subequations}
\begin{align}
\tag{2.7.1c}
u_1 &= v_1\,\gamma_1 = v_1/(1 - v_1^2)^{1/2}
  = \text{``4-velocity'' on front side},\\[4pt]
\tag{2.7.1d}
u_2 &= v_2\,\gamma_2 = v_2/(1 - v_2^2)^{1/2}
  = \text{``4-velocity'' on back side}.
\end{align}
\end{subequations}

\noindent
Similarly, energy and momentum conservation are equivalent to continuity of
energy flux and momentum flux:
%
\begin{subequations}
\begin{align}
\tag{2.7.2b}
\bigl[(\rho + p)\,v\,u\bigr]_1
  &= \bigl[(\rho + p)\,v\,u\bigr]_2,\\[4pt]
\tag{2.7.2c}
\bigl[(\rho + p)\,u^2 + p\bigr]_1
  &= \bigl[(\rho + p)\,u^2 + p\bigr]_2.
\end{align}
\end{subequations}

%%% ---------------------------------------------------------------
%%% Figure 2.7.1
%%% ---------------------------------------------------------------

\begin{figure}[htbp]
\centering
% \includegraphics[width=\columnwidth]{fig_2_7_1}
\fbox{\parbox{0.9\columnwidth}{\centering [Figure 2.7.1 placeholder]\\[6pt]
Diagram of a shock wave viewed in the local-Lorentz rest frame of the
shock front at an event~$\mathscr{P}$.  The front (``1'' side) is on the left,
and the back (``2'' side) is on the right.  Arrows indicate gas flow with
velocity~$v_1$ entering from the left and velocity~$v_2$ exiting to the right.
The $x$-axis is perpendicular to the shock surface, which lies in the
$y$-$z$ plane.  Hatching indicates the compressed gas behind the shock.}}
\caption{\textit{Figure~2.7.1.}  A shock wave viewed in the local-Lorentz
rest frame of the shock front at an event~$\mathscr{P}$.}
\label{fig:2.7.1}
\end{figure}

These junction conditions are more easily understood by writing them in a
form analogous to the Rankine--Hugoniot equations of Newtonian theory:
First take the law of baryon conservation~(2.7.2a) and turn it into equations
for the fluid 4-velocities
%
\begin{equation}
\tag{2.7.3a}
u_1 = jV_1, \qquad u_2 = jV_2.
\end{equation}
%
(See~\S2.5 for notation.)  Then take the law of momentum conservation~(2.7.2a)
and rewrite it in terms of $\mu$, $V_1$, and $p$ using the law of baryon
conservation~(2.7.2a) to obtain
%
\begin{subequations}
\begin{align}
\tag{2.7.3b}
j^2 &= -\frac{p_2 - p_1}{\mu_2\,V_2 - \mu_1\,V_1},\\[6pt]
\end{align}
\end{subequations}
%
and the relations
$j^2 = (p_2-p_1)/(\mu_2 V_2 - \mu_1 V_1)$.

Finally, use the law of baryon conservation~(2.7.2b) to rewrite the energy
equation~(2.7.3b) by $(\mu_1 V_1 + \mu_2 V_2)$, and combine with
$j = u_1/V_1 = u_2/V_2$; divide equation~(2.7.3b) by the square of the energy
equation $(\mu_2\gamma_2)^2 - (\mu_1\gamma_1)^2 = (p_1 - p_2)(\mu_1 V_1 + \mu_2 V_2)$;
then subtract this from [eq.~(2.7.2a)] to obtain
%
\begin{equation}
\tag{2.7.3c}
\mu_2^2 - \mu_1^2 = (p_2 - p_1)(\mu_1\,V_1 + \mu_2\,V_2).
\end{equation}

\emph{Equations~(2.7.3) are Taub's~(1948) junction conditions for shock waves.}
Their Newtonian limits are the standard Rankine--Hugoniot equations:
%
\begin{subequations}
\begin{align}
\tag{2.7.4a}
v_1 &= j_0\,V_{01},\qquad v_2 = j_0\,V_{02}\\[4pt]
\tag{2.7.4b}
j_0 &\equiv \frac{p_2 - p_1}{V_{01} - V_{02}}\quad
\text{(``mass flux'' of Newtonian theory)};\\[4pt]
\tag{2.7.4c}
w_2 - w_1 &= \tfrac{1}{2}(p_2 - p_1)(V_{01} + V_{02}).
\end{align}
\end{subequations}
%
The form of the Newtonian junction conditions~(2.7.4) motivates one to use $p$
and $V_0$ as one's independent thermodynamic variables when analyzing Newtonian
shocks; similarly, the form of the relativistic junction conditions~(2.7.3)
motivates one to use $p$ and~$\mu V$ as one's independent variables for
relativistic shocks.

\subsubsection*{The Rankine-Hugoniot curve}

Consider a family of shocks, each with the same thermodynamic state on the
front face (same $\mu_1$, $V_1$, $p_1$, etc.), but with different states on the
back face (different $\mu_2$, $V_2$, $p_2$, etc.).  This family of shocks is
a one-parameter family.  Thus, if one plots all back-face states
$(\mu_2 V_2,\,p_2)$ in the $\mu V$--$p$ plane, they lie on a single
curve---the ``Rankine-Hugoniot curve''---passing through the point
$(\mu_1 V_1,\,p_1)$.  (In the relativistic case this curve is also called the
``Taub adiabat'' or ``Poisson adiabat''.)

One can also plot, in the $\mu V$--$p$ plane, the Poisson adiabat (curve of
constant entropy) passing through $(\mu_1 V_1,\,p_1)$.  These two curves
typically have the relative shapes and locations shown in Figure~2.7.2.  In
particular, from the Rankine-Hugoniot equations one can derive the following
general properties of the Rankine-Hugoniot
curve:\footnote{See \citet{Thorne1973a} for derivation.  (i)~The
  Rankine-Hugoniot curve is tangent to the Poisson adiabat, and has the same
  second derivative at point~``1''.  The derivation requires the assumption
  $[\partial^2(\mu V)/\partial p^2]_{s} > 0$; a condition which is satisfied
  by all materials of astrophysical or laboratory interest; see
  \citet{LandauLifshitz1959}.  Some, but not all, of the listed properties
  remain true without this condition; see \citet{ZeldovichRaizer1966} for
  nonrelativistic details, and \citet{Thorne1973a} for relativistic details.}

\begin{enumerate}
\item[(i)] The Rankine-Hugoniot curve is tangent to the Poisson adiabat, and
  has the same second derivative at point~``1''; i.e., for weak shocks the
  increase in entropy is third-order in the pressure jump:
  %
  \begin{equation}
    s_2 - s_1 = \left[\frac{1}{12\mu T}\,
      \frac{\partial^2(\mu V)}{\partial p^2}\bigg|_{s,\,1}\right]
      (p_2 - p_1)^3 + O\!\left[(p_2-p_1)^4\right].
    \tag{2.7.5}
  \end{equation}

\item[(ii)] As the gas passes from the front of the shock to the back, its
  entropy, pressure, and chemical potential increase
  %
  \begin{equation}
    s_2 > s_1\,,\qquad p_2 > p_1\,,\qquad \mu_2 > \mu_1\,;
    \tag{2.7.6a}
  \end{equation}
  %
  while its specific volume and the product $\mu V$ decrease
  %
  \begin{equation}
    V_2 < V_1\,,\qquad \mu_2 V_2 < \mu_1 V_1\,.
    \tag{2.7.6b}
  \end{equation}
  %
  (Hence, only the ``upper branch'' of the Rankine-Hugoniot curve,
  Figure~2.7.2, is physically relevant.)

\item[(iii)] The flow on the front side is always supersonic; on the back side
  is always subsonic
  %
  \begin{equation}
    u_1/u_{S1} > 1\,,\qquad u_2/u_{S2} < 1\,;
    \tag{2.7.7}
  \end{equation}
  %
  \[
    u_S = c_S/(1 - c_S^2)^{1/2}\,.
  \]

\item[(iv)] As one moves up the Rankine-Hugoniot curve, away from point
  ``1''---i.e., as one studies a sequence of ever ``stronger''
  shocks---the following quantities increase monotonically:
  %
  \begin{equation}
  \begin{gathered}
    \text{the baryon flux across the shock, } j\,; \\
    \text{the jump in entropy across the shock, } s_2 - s_1\,; \\
    \text{the ``relativistic'' Mach number on the front side of the shock, }
      M_1 = u_1/u_{S1}\,.
  \end{gathered}
    \tag{2.7.8}
  \end{equation}
\end{enumerate}

Notice that, once the thermodynamic state on the front face of the shock
has been specified, the shock has only one free parameter.  If one fixes the
baryon flux across the shock~$j$, or the speed on the front face~$u_1$, or
the pressure on the back face~$p_2$, or any other single parameter, then all
other properties of the shock are uniquely determined.  To calculate them, one
need merely invoke the Rankine-Hugoniot equations~(2.7.4), the laws of
thermodynamics, and the equation of state of the gas.

%%% Shocks in an ideal gas %%%

\subsubsection*{Shocks in an ideal gas}

Consider, as a special but important case, a nonrelativistic ideal gas with
constant adiabatic index $\Gamma = -(\partial\ln p/\partial\ln V_0)_s$, so
that
%
\begin{equation}
  pV_0 = kT/\bar{m}\,.
  \tag{2.7.9}
\end{equation}

By a straightforward calculation (\S85 of \citealt{LandauLifshitz1959}), one
can derive the following relationships between various quantities along the
Rankine-Hugoniot curve:
%
\begin{equation}
\begin{aligned}
  \frac{V_{02}}{V_{01}} &= \frac{(\Gamma+1)\,p_1 + (\Gamma-1)\,p_2}
                               {(\Gamma-1)\,p_1 + (\Gamma+1)\,p_2}
  \;\to\; \frac{\Gamma-1}{\Gamma+1}
  &&\text{for strong shock,} \\[6pt]
  %
  \frac{T_2}{T_1} &= \frac{p_2}{p_1}\,
    \frac{(\Gamma-1)\,p_2 + (\Gamma+1)\,p_1}
         {(\Gamma+1)\,p_2 + (\Gamma-1)\,p_1}
  \;\to\; \frac{\Gamma-1}{\Gamma+1}\,\frac{p_2}{p_1}
  &&\text{for strong shock,} \\[6pt]
  %
  \hat{\jmath}_0^{\,2} &=
    \frac{(\Gamma-1)\,p_1 + (\Gamma+1)\,p_2}{2V_{01}}
  \;\to\; \frac{\Gamma+1}{2}\,\frac{p_2}{V_{01}}
  &&\text{for strong shock,} \\[6pt]
  %
  f^2 &= \tfrac{1}{2}\,V_{01}
    \bigl[(\Gamma+1)\,p_1 + (\Gamma-1)\,p_2\bigr]
  \;\to\; \tfrac{1}{2}\,(\Gamma-1)\,V_{01}\,p_2
  &&\text{for strong shock,} \\[6pt]
  %
  \hat{\jmath}_0^{\,2} &=
    \frac{V_{01}\bigl[(\Gamma+1)\,p_1 + (\Gamma-1)\,p_2\bigr]^2}
         {2(\Gamma+1)}
  \;\to\; \frac{(\Gamma-1)^2}{2(\Gamma+1)}\,V_{01}\,p_2^{\,2}
  &&\text{for strong shock.}
\end{aligned}
  \tag{2.7.10}
\end{equation}
%
Here ``strong shock'' means ``in the limit $p_2/p_1 \gg 1$''.

\begin{figure}[htbp]
\centering
\fbox{\parbox{0.9\columnwidth}{\centering [Figure 2.7.2 placeholder]\\
The Poisson adiabat $\mathscr{P}\,\mathscr{P}'$ and the Rankine-Hugoniot
curve $\mathscr{R}\,\mathscr{R}'$ plotted in the $\mu V$--$p$ plane.}}
\caption{The Poisson adiabat $\mathscr{P}\,\mathscr{P}'$ and the
  Rankine-Hugoniot curve $\mathscr{R}\,\mathscr{R}'$ plotted in the
  $\mu V$--$p$ plane.}
\label{fig:2.7.2}
\end{figure}


%%%%%%%%%%%%%%%%%%%%%%%%%%%%%%%%%%%%%%%%%%%%%%%%%%%%%%%%%%%%%%%%
%%% SECTION 2.8  TURBULENCE
%%%%%%%%%%%%%%%%%%%%%%%%%%%%%%%%%%%%%%%%%%%%%%%%%%%%%%%%%%%%%%%%

\subsection{Turbulence}
\label{sec:2.8}

In the accretion of gas onto a black hole, turbulence probably plays an
important role.  (See \S\S4 and~5.)  Unfortunately, the theory of turbulence
is in a very uncertain state.  Little is known with confidence about the
astrophysical circumstances under which turbulence should develop, or about
the strength of the turbulence; see, e.g., Chapter~3 of
\citet{LandauLifshitz1959}; also \citet{Pikelner1961}.


%%%%%%%%%%%%%%%%%%%%%%%%%%%%%%%%%%%%%%%%%%%%%%%%%%%%%%%%%%%%%%%%
%%% SECTION 2.9  RECONNECTION OF MAGNETIC FIELD LINES
%%%%%%%%%%%%%%%%%%%%%%%%%%%%%%%%%%%%%%%%%%%%%%%%%%%%%%%%%%%%%%%%

\subsection{Reconnection of Magnetic Field Lines}
\label{sec:2.9}

\emph{Basic references}\quad \S5.3 of \citet{Cowling1965};
\citet{Sonnerup1970}; \citet{Yeh1970};
\citet{VainshteinZeldovich1972}.

\subsubsection*{Basic ideas}

In flat spacetime consider a plasma that is macroscopically neutral, that has a
high electrical conductivity~$\sigma$, that is sufficiently dilute for one to
ignore its dielectric properties: $\varepsilon = \mu = 1$.  Let the plasma be
endowed with a large-scale magnetic field, and examine the evolution of that
field in a Lorentz frame where the plasma has low velocity
$|\mathbf{v}| \ll c$.  Maxwell's equations for the electromagnetic field then
read
%
\begin{align}
  \nabla \times \mathbf{E} + (1/c)(\partial\mathbf{B}/\partial t) &= 0\,,
  \tag{2.9.1} \\[4pt]
  \nabla \times \mathbf{B} - (1/c)(\partial\mathbf{E}/\partial t)
    &= 4\pi\mathbf{J}/c\,.
  \tag{2.9.2}
\end{align}

In the local rest frame of the plasma the current is proportional to the
electric field, $\mathbf{J} = \sigma\mathbf{E}$.  When transformed to the
Lorentz frame where the plasma moves with velocity~$\mathbf{v}$, this
equation says
%
\begin{equation}
  \mathbf{J} = \sigma\bigl[\mathbf{E}
    + (\mathbf{v}/c) \times \mathbf{B}\bigr].
  \tag{2.9.3}
\end{equation}

A straightforward calculation from the above equations leads to the following
law for the rate of change of magnetic field along the world lines of the
plasma:
%
\begin{equation}
  \frac{d\mathbf{B}}{d\tau}
  = (\omega + \sigma - \tfrac{2}{3}\theta\,\mathbf{1}) \cdot \mathbf{B}
    - \frac{1}{4\pi\sigma}\left(
      \frac{\partial^2\mathbf{B}}{\partial t^2}
      - c^2\,\nabla^2\mathbf{B}\right).
  \tag{2.9.4}
\end{equation}
%
Here $\mathbf{1}$ is the unit 3-tensor; $\omega$, $\sigma$, and~$\theta$ are
the rotation, shear, and expansion of the plasma
[cf.\ eqs.~(2.5.17) and~(2.5.18)]; and
%
\[
  d/d\tau = \partial/\partial t + \mathbf{v}\cdot\nabla\,.
\]

Because of the very high electrical conductivity, one can ignore the
``wave-equation'' part of eq.~(2.9.4) almost everywhere in the plasma;
and the magnetic field evolves in a ``frozen-in'' manner:
%
\begin{equation}
  d\mathbf{B}/d\tau
  = (\omega + \sigma - \tfrac{2}{3}\theta\,\mathbf{1}) \cdot \mathbf{B}\,.
  \tag{2.9.5}
\end{equation}

[Cf.\ Eq.~(2.5.20) and associated discussion.]  However, in regions of plasma
where the field has strong gradients, the ``wave-equation'' part must come
into play.

Consider, as the case of greatest importance, a magnetic field that is chaotic
with an ordered structure on some scale~$l_c$ (subscript $c$ for ``cell
size'').  At the interface between adjacent cells, the magnetic field must
reverse sign (``neutral point'' or ``neutral sheet''), and its gradients will
be very large.  Hence, at the interface, ``wave-equation'' behavior can come
into play destroying the frozen-in behavior of the field.  The result is a
reconnection of magnetic field lines, as

\begin{figure}[htbp]
\centering
\fbox{\parbox{0.9\columnwidth}{\centering [Figure 2.9.1 placeholder]\\
(a)~Top view and (b)~perspective view of reconnecting magnetic field lines
in a plasma.  The reconnection region is stippled; the magnetic field lines
are drawn solidly and the flow lines of the plasma are dotted.  The innermost
field line is in the process of reconnecting into the dashed form.
The perspective view~(b) shows the closed curve~$\mathscr{E}$ around which
one integrates to derive eq.~(2.9.6) for the EMF in the connecting region.}}
\caption{The region in a plasma where magnetic field lines are
  reconnecting (schematic picture).  In the top view, (a), the
  reconnection region is stippled; the magnetic field lines are drawn
  solidly; and the flow lines of the plasma are dotted.  The innermost
  field line is in the process of reconnecting into the dashed form.
  The perspective view, (b), shows the closed curve~$\mathscr{E}$
  around which one integrates to derive eq.~(2.9.6) for the EMF in the
  connecting region.}
\label{fig:2.9.1}
\end{figure}

depicted in Figure~2.9.1.  This reconnection gradually converts the two adjacent
cells into a single cell, reducing the overall strength of the magnetic field.

In situations of interest (\S\S4.4 and 5.2), the decrease of field strength due
to reconnection will be counterbalanced by an increase in strength due to
compression or shear of the plasma.

The reconnection process creates very strong electric fields.  In particular,
a simple application of Faraday's law
%
\begin{equation}
\oint_{\mathscr{C}} \mathbf{E}\cdot d\mathbf{l}
  = -\frac{d}{d\tau}\!\int \mathbf{B}\cdot d\mathbf{A},
\end{equation}
%
with the curve~$\mathscr{C}$ as shown in Figure~2.9.1b, shows that the EMF built up
in the reconnecting region is
%
\begin{equation}
\Delta\phi_e
  = (\text{rate of reconnection of magnetic flux})
  \equiv d\Psi/dt.
\tag{2.9.6}
\end{equation}

In situations of interest to us [\S5.12; Lynden-Bell~(1969)\cite{LyndenBell1969}],
these EMF's can be far greater than $10^{10}$~volts.  If the density of the plasma is
sufficiently low to permit long mean-free paths for charged particles, and if the
reconnecting region is sufficiently straight, then these EMF's can accelerate
particles to ultra-relativistic energies; and the particles can then radiate intense
synchrotron radiation.  In this manner regions of reconnecting field lines may be
sources of flares.
